\chapter{Index}\label{ch:index}

Les index sont produits automatiquement à partir du texte transmis. Ils renvoient aux numéros des notices, non aux pages. L’astérisque signale une forme dont une lettre au moins est restituée ; dans l’index des consuls, le cercle (°) signale une lecture partielle du nom et l’astérisque un nom restitué ou non lu. Les formes sont données telles qu’elles sont transmises, sans lemmatisation : un même mot peut apparaître sous plusieurs formes.

\section{Index des noms de personnes et de lieux}

Noms à initiale capitale, hors noms consulaires de l’année du texte (chapitre~\ref{ch:noms}).

% généré par scripts/04_generated.py
\begin{multicols}{3}\footnotesize\setlength{\parindent}{0pt}\raggedright
\hangindent=1em Abita 330\par
\hangindent=1em Abudantia 230\par
\hangindent=1em Abundanti 448\par
\hangindent=1em Accill 339\par
\hangindent=1em Acone 32*\par
\hangindent=1em Adelfius 36\par
\hangindent=1em Adeoda 378*\par
\hangindent=1em Adeodatae 297\par
\hangindent=1em Afrodisi 39*\par
\hangindent=1em Agape 207\par
\hangindent=1em Agapeni 99\par
\hangindent=1em Agapenis 134\par
\hangindent=1em Agrige 147\par
\hangindent=1em Agrippina 133\par
\hangindent=1em Albanensis 378\par
\hangindent=1em Alexander 115\par
\hangindent=1em Alexandra 333*, 404\par
\hangindent=1em Alexandrae 367\par
\hangindent=1em Alexandriae 304\par
\hangindent=1em Alicensi 164\par
\hangindent=1em Amachius 23\par
\hangindent=1em Amantiaes 383\par
\hangindent=1em Amantio 166\par
\hangindent=1em Ampliato 399\par
\hangindent=1em Anastasa 197\par
\hangindent=1em Anastasi 195\par
\hangindent=1em Anastasio 43, 321\par
\hangindent=1em Anicitionis 37\par
\hangindent=1em Antus 107\par
\hangindent=1em Apolinus 102\par
\hangindent=1em Apro 185\par
\hangindent=1em Apronianae 279\par
\hangindent=1em Aqundino 55\par
\hangindent=1em Aradius 26\par
\hangindent=1em Aristo 308\par
\hangindent=1em Armeniacum 275\par
\hangindent=1em Artemia 137\par
\hangindent=1em Artorius 23\par
\hangindent=1em Asello 100\par
\hangindent=1em Asterius 95\par
\hangindent=1em Astuto 280*\par
\hangindent=1em Aufidia 449\par
\hangindent=1em Aufinianus 391\par
\hangindent=1em Augendo 146\par
\hangindent=1em Aulo 339\par
\hangindent=1em Aur 370\par
\hangindent=1em Aureli 19*\par
\hangindent=1em Aureliae 6, 23*\par
\hangindent=1em Aurelio 111\par
\hangindent=1em Aurelius 6*, 19, 229\par
\hangindent=1em Ausustinae 84\par
\hangindent=1em Avi 24\par
\hangindent=1em Balbilla 72\par
\hangindent=1em Barbarus 478\par
\hangindent=1em Basili 378\par
\hangindent=1em Beatus 381\par
\hangindent=1em Benata 407\par
\hangindent=1em Benedictus 268*, 441\par
\hangindent=1em Benenato 110\par
\hangindent=1em Benerio 402\par
\hangindent=1em Benerosa 323\par
\hangindent=1em Bictora 4\par
\hangindent=1em Bincentiae 369\par
\hangindent=1em Bonae 92*\par
\hangindent=1em Bonifae 405\par
\hangindent=1em Bonifatia 84*\par
\hangindent=1em Bonifatiae 140\par
\hangindent=1em Bonoso 16\par
\hangindent=1em Caelestinae 399\par
\hangindent=1em Caelia 204\par
\hangindent=1em Caiani 22\par
\hangindent=1em Calchedonia 179\par
\hangindent=1em Calleroe 193\par
\hangindent=1em Calliste 81, 181*\par
\hangindent=1em Callistratus 181*\par
\hangindent=1em Cappadocem 275\par
\hangindent=1em Carabasinus 57\par
\hangindent=1em Casarius 19\par
\hangindent=1em Cassia 15\par
\hangindent=1em Catafroniaes 134*\par
\hangindent=1em Cecili 102\par
\hangindent=1em Celerine 231\par
\hangindent=1em Celerini 44*\par
\hangindent=1em Celerinus 347\par
\hangindent=1em Celestinus 297\par
\hangindent=1em Cerealis 319\par
\hangindent=1em Cheretio 17\par
\hangindent=1em Cheretius 17\par
\hangindent=1em Cinnamius 202\par
\hangindent=1em Clodius 36\par
\hangindent=1em Concordiaes 424\par
\hangindent=1em Concordius 59, 141\par
\hangindent=1em Constant 138\par
\hangindent=1em Constantiae 455\par
\hangindent=1em Constantius 147\par
\hangindent=1em Cornificiae 77\par
\hangindent=1em Crescens 125\par
\hangindent=1em Crescentia 425\par
\hangindent=1em Crescentiaes 417\par
\hangindent=1em Crescentii 378*\par
\hangindent=1em Crescentina 440\par
\hangindent=1em Criscent 331\par
\hangindent=1em Criscentinae 132\par
\hangindent=1em Damasus 244\par
\hangindent=1em Denudatae 442\par
\hangindent=1em Didyme 179\par
\hangindent=1em Diocletiano 64*\par
\hangindent=1em Dionisus 300\par
\hangindent=1em Domitio 157\par
\hangindent=1em Donata 364, 373\par
\hangindent=1em Donatus 27\par
\hangindent=1em Dulciti 134\par
\hangindent=1em Egnatia 330\par
\hangindent=1em Eleuteriae 73\par
\hangindent=1em Elpidius 206\par
\hangindent=1em Emerita 339\par
\hangindent=1em Epinicenis 418\par
\hangindent=1em Erenes 78\par
\hangindent=1em Eros 387\par
\hangindent=1em Esuviae 26\par
\hangindent=1em Eudaemon 62*\par
\hangindent=1em Eudoxiae 265\par
\hangindent=1em Euodius 33\par
\hangindent=1em Eusebi 309\par
\hangindent=1em Eusebius 191\par
\hangindent=1em Eustolius 213\par
\hangindent=1em Eutyches 68\par
\hangindent=1em Eutychius 352, 352*\par
\hangindent=1em Exsuperantius 156\par
\hangindent=1em Exsuperia 49\par
\hangindent=1em Extricatus 171*\par
\hangindent=1em Exuperantius 245*\par
\hangindent=1em Fasciola 421*\par
\hangindent=1em Fasciolae 202\par
\hangindent=1em Faustinae 134*\par
\hangindent=1em Faustinus 205\par
\hangindent=1em Feliciae 351\par
\hangindent=1em Felicis 102\par
\hangindent=1em Felicissima 34\par
\hangindent=1em Felicitas 301, 346*, 470\par
\hangindent=1em Felicitatis 417*\par
\hangindent=1em Felix 22, 81, 133, 229\par
\hangindent=1em Festo 117\par
\hangindent=1em Filippo 321\par
\hangindent=1em Filonetis 250\par
\hangindent=1em Flavi 98\par
\hangindent=1em Flavia 72, 440\par
\hangindent=1em Flaviae 221\par
\hangindent=1em Flavio 157, 321\par
\hangindent=1em Flavius 19, 20, 57*, 59, 149, 217\par
\hangindent=1em Florentia 157, 472\par
\hangindent=1em Flori 244\par
\hangindent=1em Floridus 266\par
\hangindent=1em Fortunius 204\par
\hangindent=1em Gaianus 20*, 94\par
\hangindent=1em Galatea 151\par
\hangindent=1em Galla 371\par
\hangindent=1em Gaudentia 451\par
\hangindent=1em Gaudentius 34, 310, 372\par
\hangindent=1em Gelasio 131\par
\hangindent=1em Germana 208*\par
\hangindent=1em Gerontio 211\par
\hangindent=1em Gregori 134*\par
\hangindent=1em Heraclius 347\par
\hangindent=1em Herculius 159*\par
\hangindent=1em Hermogenia 10*\par
\hangindent=1em Hermoginiano 371*\par
\hangindent=1em Hesperius 73*\par
\hangindent=1em Hilara 116\par
\hangindent=1em Hilarino 117, 410*\par
\hangindent=1em Hilaritati 324, 399\par
\hangindent=1em Hilaro 320\par
\hangindent=1em Hispanus 296\par
\hangindent=1em Hyppolytus 98\par
\hangindent=1em Ianuaria 9, 118*\par
\hangindent=1em Ianuariae 26\par
\hangindent=1em Ianus 43*\par
\hangindent=1em Innocentia 85*\par
\hangindent=1em Innocentius 67*, 179\par
\hangindent=1em Ioannti 261\par
\hangindent=1em Iosimus 365\par
\hangindent=1em Ioviano 247\par
\hangindent=1em Iovina 138*\par
\hangindent=1em Iuliae 221\par
\hangindent=1em Iulius 67*\par
\hangindent=1em Iuniani 4\par
\hangindent=1em Iuniori 191\par
\hangindent=1em Iustina 45\par
\hangindent=1em Iustinus 346, 350*\par
\hangindent=1em Iustus 395\par
\hangindent=1em Karisiae 46\par
\hangindent=1em Lacurius 65*\par
\hangindent=1em Lampadius 247\par
\hangindent=1em Laur 217*\par
\hangindent=1em Laurenti 89, 237, 284*\par
\hangindent=1em Laurentia 233, 257*, 303, 312*, 339, 432*\par
\hangindent=1em Laurentio 92\par
\hangindent=1em Laurentius 428*, 463\par
\hangindent=1em Lea 19, 19*, 97, 272*, 347\par
\hangindent=1em Legitima 47\par
\hangindent=1em Leo 7, 60, 194, 295, 333, 387\par
\hangindent=1em Leontia 125, 460\par
\hangindent=1em Leontius 276\par
\hangindent=1em Leoparda 203, 349\par
\hangindent=1em Leopardes 102\par
\hangindent=1em Leopardo 135\par
\hangindent=1em Leopardus 269\par
\hangindent=1em Lepus 430*\par
\hangindent=1em Libera 172\par
\hangindent=1em Ligitimus 100\par
\hangindent=1em Limenio 378*\par
\hangindent=1em Limenius 378*\par
\hangindent=1em Limfiria 185\par
\hangindent=1em Lucceia 306\par
\hangindent=1em Lucida 318\par
\hangindent=1em Lucifero 268\par
\hangindent=1em Lucio 433\par
\hangindent=1em Marasius 41\par
\hangindent=1em Marcelli 378\par
\hangindent=1em Marcellinae 194\par
\hangindent=1em Marcellinus 55\par
\hangindent=1em Marcena 152\par
\hangindent=1em Marciana 454*\par
\hangindent=1em Marcus 46, 197\par
\hangindent=1em Mariniano 59\par
\hangindent=1em Martia 275\par
\hangindent=1em Martielis 429\par
\hangindent=1em Martinum 454\par
\hangindent=1em Martyrius 139*\par
\hangindent=1em Maura 130\par
\hangindent=1em Maxima 419\par
\hangindent=1em Maximae 144\par
\hangindent=1em Maximi 411\par
\hangindent=1em Maximiano 396\par
\hangindent=1em Maximimasiae 320\par
\hangindent=1em Maximo 72\par
\hangindent=1em Melissus 26\par
\hangindent=1em Mercurius 422*\par
\hangindent=1em Muscula 151\par
\hangindent=1em Nae 221\par
\hangindent=1em Niceforidi 167\par
\hangindent=1em Nicetius 296\par
\hangindent=1em Nisus 311\par
\hangindent=1em Nunita 306\par
\hangindent=1em Olibioni 260\par
\hangindent=1em Opas 202\par
\hangindent=1em Pacactianus 73*\par
\hangindent=1em Pacatiano 73\par
\hangindent=1em Palladio 344*\par
\hangindent=1em Palma 382*\par
\hangindent=1em Pannonia 306\par
\hangindent=1em Panthagate 382*\par
\hangindent=1em Pasc 467\par
\hangindent=1em Pascasio 400\par
\hangindent=1em Pascasius 235\par
\hangindent=1em Pascentius 421\par
\hangindent=1em Paterniae 223\par
\hangindent=1em Paula 160*\par
\hangindent=1em Pauli 218, 321, 387\par
\hangindent=1em Paulinus 6\par
\hangindent=1em Paulus 398\par
\hangindent=1em Pelagia 246\par
\hangindent=1em Pelecrinus 56\par
\hangindent=1em Perseberatia 299\par
\hangindent=1em Petiliano 329*, 336\par
\hangindent=1em Petri 154, 387*\par
\hangindent=1em Petrio 68\par
\hangindent=1em Petronia 416\par
\hangindent=1em Petroniae 184\par
\hangindent=1em Petroniaes 183\par
\hangindent=1em Petronio 144, 145, 149*\par
\hangindent=1em Petronius 374\par
\hangindent=1em Petrus 197\par
\hangindent=1em Petulum 439\par
\hangindent=1em Pisonis 15\par
\hangindent=1em Pitzinnina 346\par
\hangindent=1em Porfurius 204\par
\hangindent=1em Praetextatus 188*\par
\hangindent=1em Praetioso 458\par
\hangindent=1em Preiecticio 424\par
\hangindent=1em Preniste 19\par
\hangindent=1em Prima 54\par
\hangindent=1em Primi 10*\par
\hangindent=1em Primia 254\par
\hangindent=1em Primo 244\par
\hangindent=1em Primus 10\par
\hangindent=1em Priscus 132\par
\hangindent=1em Privata 38\par
\hangindent=1em Probi 143\par
\hangindent=1em Probinianus 445*\par
\hangindent=1em Proclino 387\par
\hangindent=1em Proiectae 244\par
\hangindent=1em Pudentiana 269\par
\hangindent=1em Quiminati 210\par
\hangindent=1em Quinta 50\par
\hangindent=1em Quinti 350\par
\hangindent=1em Quintiano 47\par
\hangindent=1em Quintiliano 366\par
\hangindent=1em Quintilianus 474\par
\hangindent=1em Quinto 339, 340\par
\hangindent=1em Quiriace 298\par
\hangindent=1em Quirillus 275\par
\hangindent=1em Quodvuldeus 112*\par
\hangindent=1em Quodvultdeus 388\par
\hangindent=1em Quriaceti 319\par
\hangindent=1em Rapetiga 296\par
\hangindent=1em Refrigerio 75\par
\hangindent=1em Remigia 428\par
\hangindent=1em Rerinta 390\par
\hangindent=1em Reveccae 409\par
\hangindent=1em Romanillae 430*\par
\hangindent=1em Romano 145\par
\hangindent=1em Romanus 22\par
\hangindent=1em Romuliano 155\par
\hangindent=1em Roscia 179\par
\hangindent=1em Rufinae 281*\par
\hangindent=1em Rufino 29\par
\hangindent=1em Rufinus 142*\par
\hangindent=1em Sabbatius 22, 386\par
\hangindent=1em Sabina 200\par
\hangindent=1em Sabiniane 286\par
\hangindent=1em Sambatiae 339*\par
\hangindent=1em Santulus 19*\par
\hangindent=1em Satullo 111\par
\hangindent=1em Saturnino 462*\par
\hangindent=1em Seberane 124\par
\hangindent=1em Secundae 70\par
\hangindent=1em Secundi 48\par
\hangindent=1em Secundianaes 66\par
\hangindent=1em Secundinus 136\par
\hangindent=1em Serbulus 235\par
\hangindent=1em Setema 387\par
\hangindent=1em Severa 17\par
\hangindent=1em Severae 310\par
\hangindent=1em Severianus 400\par
\hangindent=1em Severini 98*\par
\hangindent=1em Sextilius 343*\par
\hangindent=1em Sil 335\par
\hangindent=1em Silvana 56\par
\hangindent=1em Simplicius 179\par
\hangindent=1em Sirica 459\par
\hangindent=1em Sirici 387*\par
\hangindent=1em Siricius 321\par
\hangindent=1em Socratiaes 48\par
\hangindent=1em Sosoro 119\par
\hangindent=1em Soteres 247\par
\hangindent=1em Sozomene 97\par
\hangindent=1em Statilia 42\par
\hangindent=1em Staturina 386\par
\hangindent=1em Stercoria 108*\par
\hangindent=1em Successae 35\par
\hangindent=1em Susanna 330\par
\hangindent=1em Sussanna 406\par
\hangindent=1em Taeodore 6*\par
\hangindent=1em Tertulla 61\par
\hangindent=1em Theodora 80\par
\hangindent=1em Theodosius 452\par
\hangindent=1em Theodotus 5*\par
\hangindent=1em Theofila 110\par
\hangindent=1em Ticira 255\par
\hangindent=1em Tigris 42\par
\hangindent=1em Tigriti 353\par
\hangindent=1em Titiana 207\par
\hangindent=1em Titiani 418*\par
\hangindent=1em Urania 241*\par
\hangindent=1em Urbicae 284\par
\hangindent=1em Ursa 411\par
\hangindent=1em Ursi 169\par
\hangindent=1em Ursini 378\par
\hangindent=1em Urso 387*, 478\par
\hangindent=1em Ursus 124, 184, 443\par
\hangindent=1em Valeriae 109\par
\hangindent=1em Valerius 262\par
\hangindent=1em Venerandae 472\par
\hangindent=1em Vicentia 381\par
\hangindent=1em Victor 217, 373, 420, 466\par
\hangindent=1em Victori 233\par
\hangindent=1em Victoria 258*, 378*\par
\hangindent=1em Victoriae 361\par
\hangindent=1em Victoriaes 114\par
\hangindent=1em Victorinae 19*, 127*\par
\hangindent=1em Victorini 209\par
\hangindent=1em Victorino 357\par
\hangindent=1em Victorio 412\par
\hangindent=1em Vincentiae 271\par
\hangindent=1em Vincentius 105, 446\par
\hangindent=1em Viricunda 400\par
\hangindent=1em Vitalinae 376\par
\hangindent=1em Vitalini 186*\par
\hangindent=1em Viventi 306\par
\hangindent=1em Volusiana 449\par
\hangindent=1em Yperechiam 275\par
\hangindent=1em Zita 340\par
\hangindent=1em Zosimi 424\par
\hangindent=1em Zosimus 106, 120\par\end{multicols}


\section{Index des consuls}

% généré par scripts/04_generated.py
\begin{multicols}{2}\footnotesize\setlength{\parindent}{0pt}\raggedright
\hangindent=1em Acindynus : 1, 2°, 3°, 55*\par
\hangindent=1em Albinus : 11*, 12, 13*, 14*, 15*, 16, 17, 18, 19\par
\hangindent=1em Amantius : 11, 12, 13, 14, 15, 16, 17, 18°, 19\par
\hangindent=1em Anicii : 369, 371, 372, 373, 377, 379, 381, 383\par
\hangindent=1em Antonius : 228°, 229, 230, 231, 232*, 233, 234*, 235, 236*, 237°, 238*, 239*\par
\hangindent=1em Arbitio : 41, 42, 43, 44, 45°\par
\hangindent=1em Arcadius : 276, 277, 279, 280, 281*, 343, 345*, 346, 347, 363*, 367, 373°, 385*, 386, 387, 389*, 390, 391, 392, 393*, 394°, 395, 431, 434, 460*, 461, 462, 463, 464*, 465*, 466, 467*, 468*\par
\hangindent=1em Arintheus : 156, 157, 158*, 159°, 161*, 164*, 165*\par
\hangindent=1em Atticus : 396*, 397°, 398*, 399, 400, 401*, 402°, 403, 404, 405, 406, 407*, 408°, 409, 410, 412*, 413, 414*, 415*, 416\par
\hangindent=1em Ausonius : 210, 211\par
\hangindent=1em Bauto : 276, 279, 280*, 281\par
\hangindent=1em Caesarius : 396, 397, 398, 399, 400, 401°, 402*, 403, 404, 405, 406, 407, 408, 409, 410°, 411, 413*, 414°, 415*, 416\par
\hangindent=1em Catullinus : 30*, 31°, 32°, 33, 34\par
\hangindent=1em Cerealis : 49, 50, 52*, 53\par
\hangindent=1em Clearchus : 262, 263°, 264°, 265, 266, 269, 270*, 271, 272°, 274°, 275\par
\hangindent=1em Constans : 5°, 6*, 138, 455\par
\hangindent=1em Constantius : 5°, 6°, 46°, 47*, 48*, 56, 57, 58, 138, 455\par
\hangindent=1em Dagalaifus : 98*, 99, 100, 101°, 103°, 104°, 105, 109\par
\hangindent=1em Datianus : 49, 50, 52*, 53*\par
\hangindent=1em Decentius : 12*, 37, 38, 39*, 40°\par
\hangindent=1em Equitius : 172, 173°, 174*, 177°, 178°, 179, 181*, 182*, 183, 184, 185, 186, 187°, 188*\par
\hangindent=1em Eucherius : 220°, 221, 223, 224°, 225, 227°\par
\hangindent=1em Eugenius : 351, 352, 353, 354, 355, 356°, 357, 358, 359, 361\par
\hangindent=1em Euodius : 282, 283°, 284°, 285°, 286°, 287*, 288, 289°, 290*\par
\hangindent=1em Eusebius : 20*, 21°, 54\par
\hangindent=1em Eutropius : 291, 292, 293, 294°\par
\hangindent=1em Eutychianus : 417, 418, 419, 420°, 422*, 423, 424, 425°, 428°, 431\par
\hangindent=1em Florentius : 59, 60, 61\par
\hangindent=1em Fravitta : 448, 449, 450°, 451, 452*, 454, 455*, 456, 458\par
\hangindent=1em Gratianus : 99, 100, 101*, 103*, 104*, 105, 109, 143, 144, 145, 146, 147, 148*, 149, 150, 151, 152*, 153*, 154, 172, 173*, 174, 175, 176°, 177*, 178*, 179°, 181, 182°, 183, 184, 185, 186, 187°, 188, 195, 196*, 197, 198*, 199*, 200, 201*, 202, 203*, 204, 205, 213, 215, 216°, 217, 218\par
\hangindent=1em Honorius : 282, 283, 284*, 285*, 286, 287, 288, 289, 290*, 367°, 373*, 385, 386, 387*, 389, 390*, 391, 392°, 393*, 394*, 395, 417, 418*, 419, 420, 421, 422°, 423*, 424, 425, 426*, 427, 428°, 460, 461°, 463, 464, 465*, 467, 468*\par
\hangindent=1em Hypatius : 54\par
\hangindent=1em Iovianus : 83*, 85, 86°, 87, 88, 89, 90°, 91\par
\hangindent=1em Iovinus : 106, 107, 108*, 110, 111, 112, 113°, 114, 115*, 116*, 117, 118, 119, 120, 121*, 122*, 124\par
\hangindent=1em Iulianus : 46°, 47, 48*, 56, 57, 58, 73, 74*, 75, 76*, 77, 78, 79, 80, 81, 82°\par
\hangindent=1em Leontius : 10*\par
\hangindent=1em Limenius : 30°, 31*, 32°, 33, 34\par
\hangindent=1em Lollianus : 41, 42, 44*, 45\par
\hangindent=1em Lupicinus : 106, 107*, 108, 110, 111, 112*, 113*, 114, 115, 116, 117, 118°, 119, 120, 121°, 122, 124\par
\hangindent=1em Magnentius : 40\par
\hangindent=1em Magnus Maximus : 295, 296, 297, 300, 302°, 311\par
\hangindent=1em Mamertinus : 62, 63°, 64, 65, 66, 67°, 68, 69°, 70*, 71*, 72\par
\hangindent=1em Marcellinus : 4\par
\hangindent=1em Merobaudes : 195, 196, 197°, 198°, 199°, 200, 201, 202, 203*, 204, 205, 236, 241°, 242*, 243*, 244, 245*, 246*, 247, 248°, 249, 250, 251, 252*, 253*, 254°, 256*, 257, 258*, 259, 261*, 278°, 298, 299, 303\par
\hangindent=1em Modestus : 155, 156, 158°, 159*, 160, 161*, 162°, 163*, 164°, 165°\par
\hangindent=1em Neoterius : 314*, 317, 319, 320, 321, 322, 323, 324, 325°, 327°\par
\hangindent=1em Nevitta : 62°, 63, 64*, 65*, 66, 67, 68, 70*, 71°, 72\par
\hangindent=1em Nicomachus Flavianus : 362*, 364, 365, 366\par
\hangindent=1em Nigrinianus : 35, 126\par
\hangindent=1em Olybrius : 210, 211, 212, 368, 370, 371, 372, 373, 374*, 375*, 376, 378*, 379, 380*, 381, 382, 383\par
\hangindent=1em Paulus : 12°, 37, 38, 39\par
\hangindent=1em Philippus : 22, 23, 24, 25°, 26, 27, 28*, 29\par
\hangindent=1em Placidus : 7, 9\par
\hangindent=1em Probinus : 4, 368, 370*, 371, 372, 373*, 374, 375°, 376, 378°, 379*, 380, 381, 382°, 383, 384*\par
\hangindent=1em Probus : 143, 144, 145, 146, 147, 148°, 149°, 150, 151, 152, 153, 154\par
\hangindent=1em Proculus : 1*, 2°, 3*, 55\par
\hangindent=1em Promotus : 304, 305°, 306, 307, 308, 309, 310, 312*, 313, 315, 316°, 317*, 318°\par
\hangindent=1em Richomeres : 260, 262*, 263*, 264*, 265°, 266*, 269, 270, 272, 273*, 274*, 275\par
\hangindent=1em Romulus : 7, 8, 9\par
\hangindent=1em Rufinus : 15, 20, 21*, 300, 343°, 344, 345, 346, 347\par
\hangindent=1em Rumoridus : 469, 470, 471°, 472*, 473°, 474*, 475°, 476, 477°, 478\par
\hangindent=1em Salia : 23, 24*, 25*, 26, 27, 28, 29\par
\hangindent=1em Sallustius : 10*, 73, 74°, 75, 76, 77, 78°, 79, 80, 81, 82*\par
\hangindent=1em Saturninus : 236*, 241*, 242°, 243°, 244, 245, 246, 247*, 248*, 250, 251°, 252°, 253°, 254*, 255, 256*, 257°, 258, 259*, 261°, 278°\par
\hangindent=1em Sergius : 126\par
\hangindent=1em Severinus : 98\par
\hangindent=1em Stilicho : 439, 440, 441, 442°, 443, 444, 446\par
\hangindent=1em Syagrius : 219*, 221°, 222*, 223, 224*, 225, 226, 228*, 229*, 230*, 231, 232°, 233, 234, 236, 237°, 238°, 239°, 240°, 328*\par
\hangindent=1em Symmachus : 329, 330, 331*, 332*, 333, 335*, 336, 337°, 338, 339, 342°, 350\par
\hangindent=1em Tatianus : 329, 330°, 331°, 332, 333, 334°, 335, 336, 337, 338*, 339, 340, 341*, 342*, 350\par
\hangindent=1em Taurus : 59, 60, 61\par
\hangindent=1em Theodorus : 429, 430, 432*, 433, 436°, 437°, 438*\par
\hangindent=1em Theodosius : 213*, 215, 216*, 217, 218*, 298, 299, 303, 351, 352, 353, 354°, 355, 357, 358*, 359, 360°, 361*, 466, 469, 470, 471*, 472, 473*, 474, 475, 476*, 477*\par
\hangindent=1em Timasius : 304, 305*, 306, 308, 310, 312°, 313°, 315*, 316*, 317, 318*\par
\hangindent=1em Valens : 92, 93°, 94, 95, 96, 97, 117, 125, 126, 128, 129*, 130, 131, 136, 138*, 139, 140, 142, 166, 167°, 168, 169, 170, 171, 189, 190°, 191, 192°, 193, 194, 206, 207, 208, 209\par
\hangindent=1em Valentinianus : 92, 93*, 94, 95, 96, 97, 102, 117, 125, 126, 127, 128, 129*, 130, 131°, 132, 133, 134°, 135°, 136, 137*, 138*, 139°, 140, 141, 142, 166, 167°, 168, 169, 170*, 171, 189°, 190°, 191, 192°, 193, 194*, 206, 207, 208, 209, 291*, 292, 293*, 294°, 301, 314*, 317°, 319, 320, 321, 322, 323, 324, 325*, 327*\par
\hangindent=1em Varronianus : 83, 85*, 86*, 87*, 88*, 89*, 90\par
\hangindent=1em Victor : 132, 134, 135, 412\par
\hangindent=1em Vincentius : 448, 449, 450, 451, 452*, 453, 456*, 458, 459\par\end{multicols}


\section{Index du calendrier}

Dates du jour converties dans le calendrier julien, avec l’événement daté (dépo. : déposition ; déc. : décès ; nais. : naissance ; ach. : achat ; indét. : événement indéterminé).

% généré par scripts/04_generated.py
\begin{multicols}{3}\footnotesize\setlength{\parindent}{0pt}\raggedright
\hangindent=1em 1 janvier : 466 (dépo.)\par
\hangindent=1em 8 janvier : 425 (indét.)\par
\hangindent=1em 9 janvier : 350 (dépo.)\par
\hangindent=1em 10 janvier : 188 (dépo.), 298 (dépo.)\par
\hangindent=1em 12 janvier : 446 (dépo.)\par
\hangindent=1em 13 janvier : 33 (indét.), 378 (dépo.)\par
\hangindent=1em 17 janvier : 297 (dépo.)\par
\hangindent=1em 22 janvier : 301 (dépo.), 376 (dépo.)\par
\hangindent=1em 26 janvier : 102 (indét.)\par
\hangindent=1em 28 janvier : 72 (indét.)\par
\hangindent=1em 29 janvier : 469 (dépo.)\par
\hangindent=1em 30 janvier : 28 (déc.)\par
\hangindent=1em 1 février : 127 (dépo.), 134 (indét.), 337 (indét.)\par
\hangindent=1em 2 février : 38 (dépo.), 76 (dépo.)\par
\hangindent=1em 3 février : 147 (indét.)\par
\hangindent=1em 7 février : 200 (indét.)\par
\hangindent=1em 8 février : 48 (dépo.), 311 (dépo.)\par
\hangindent=1em 13 février : 400 (indét.)\par
\hangindent=1em 16 février : 109 (dépo.)\par
\hangindent=1em 19 février : 89 (indét.), 383 (dépo.)\par
\hangindent=1em 20 février : 202 (dépo.)\par
\hangindent=1em 21 février : 20 (indét.)\par
\hangindent=1em 22 février : 20 (dépo.)\par
\hangindent=1em 24 février : 185 (dépo.)\par
\hangindent=1em 29 février : 46 (dépo.)\par
\hangindent=1em 4 mars : 303 (dépo.)\par
\hangindent=1em 6 mars : 131 (dépo.)\par
\hangindent=1em 9 mars : 144 (déc.)\par
\hangindent=1em 10 mars : 275 (indét.)\par
\hangindent=1em 11 mars : 134 (indét.)\par
\hangindent=1em 18 mars : 169 (dépo.)\par
\hangindent=1em 21 mars : 318 (dépo.)\par
\hangindent=1em 23 mars : 197 (déc.), 386 (dépo.)\par
\hangindent=1em 27 mars : 136 (déc.)\par
\hangindent=1em 30 mars : 26 (dépo.), 85 (dépo.)\par
\hangindent=1em 2 avril : 310 (dépo.)\par
\hangindent=1em 12 avril : 4 (dépo.)\par
\hangindent=1em 14 avril : 478 (indét.)\par
\hangindent=1em 16 avril : 249 (indét.)\par
\hangindent=1em 19 avril : 134 (indét.)\par
\hangindent=1em 23 avril : 58 (indét.), 111 (dépo.), 412 (dépo.)\par
\hangindent=1em 24 avril : 372 (dépo.)\par
\hangindent=1em 27 avril : 441 (dépo.)\par
\hangindent=1em 29 avril : 255 (dépo.)\par
\hangindent=1em 30 avril : 143 (dépo.)\par
\hangindent=1em 1 mai : 29 (dépo.)\par
\hangindent=1em 3 mai : 179 (dépo.)\par
\hangindent=1em 5 mai : 172 (dépo.)\par
\hangindent=1em 7 mai : 57 (dépo.), 259 (indét.)\par
\hangindent=1em 9 mai : 156 (dépo.)\par
\hangindent=1em 11 mai : 329 (dépo.), 336 (dépo.)\par
\hangindent=1em 13 mai : 75 (dépo.), 106 (indét.), 120 (déc.), 295 (indét.), 365 (dépo.), 381 (déc.)\par
\hangindent=1em 14 mai : 295 (dépo.), 402 (dépo.)\par
\hangindent=1em 18 mai : 124 (dépo.)\par
\hangindent=1em 20 mai : 67 (dépo.)\par
\hangindent=1em 21 mai : 381 (déc.)\par
\hangindent=1em 23 mai : 309 (déc.)\par
\hangindent=1em 25 mai : 399 (dépo.)\par
\hangindent=1em 28 mai : 94 (indét.)\par
\hangindent=1em 29 mai : 73 (dépo.)\par
\hangindent=1em 30 mai : 357 (dépo.)\par
\hangindent=1em 1 juin : 97 (dépo.), 460 (dépo.)\par
\hangindent=1em 4 juin : 308 (déc.)\par
\hangindent=1em 5 juin : 50 (dépo.), 122 (dépo.)\par
\hangindent=1em 9 juin : 54 (indét.)\par
\hangindent=1em 12 juin : 458 (indét.)\par
\hangindent=1em 15 juin : 15 (nais.)\par
\hangindent=1em 17 juin : 198 (indét.)\par
\hangindent=1em 20 juin : 476 (indét.)\par
\hangindent=1em 21 juin : 55 (nais.), 105 (dépo.)\par
\hangindent=1em 22 juin : 95 (indét.), 276 (indét.)\par
\hangindent=1em 28 juin : 405 (dépo.)\par
\hangindent=1em 29 juin : 90 (dépo.), 416 (dépo.)\par
\hangindent=1em 1 juillet : 42 (indét.), 470 (dépo.)\par
\hangindent=1em 3 juillet : 282 (dépo.), 288 (dépo.)\par
\hangindent=1em 4 juillet : 399 (dépo.)\par
\hangindent=1em 5 juillet : 7 (indét.)\par
\hangindent=1em 6 juillet : 449 (indét.)\par
\hangindent=1em 9 juillet : 34 (dépo.), 467 (dépo.)\par
\hangindent=1em 10 juillet : 19 (déc.)\par
\hangindent=1em 11 juillet : 262 (indét.)\par
\hangindent=1em 12 juillet : 27 (déc.)\par
\hangindent=1em 14 juillet : 343 (dépo.)\par
\hangindent=1em 15 juillet : 70 (déc.), 211 (déc.), 443 (dépo.)\par
\hangindent=1em 16 juillet : 211 (dépo.)\par
\hangindent=1em 18 juillet : 63 (dépo.), 151 (dépo.), 320 (dépo.)\par
\hangindent=1em 19 juillet : 12 (dépo.)\par
\hangindent=1em 21 juillet : 319 (dépo.)\par
\hangindent=1em 22 juillet : 49 (dépo.), 59 (dépo.)\par
\hangindent=1em 23 juillet : 213 (dépo.), 390 (dépo.)\par
\hangindent=1em 24 juillet : 300 (déc.)\par
\hangindent=1em 25 juillet : 273 (dépo.)\par
\hangindent=1em 26 juillet : 406 (ach.)\par
\hangindent=1em 2 août : 22 (dépo.)\par
\hangindent=1em 3 août : 369 (indét.)\par
\hangindent=1em 6 août : 459 (dépo.)\par
\hangindent=1em 11 août : 126 (indét.), 366 (dépo.)\par
\hangindent=1em 14 août : 233 (dépo.)\par
\hangindent=1em 15 août : 35 (dépo.), 304 (dépo.)\par
\hangindent=1em 16 août : 226 (dépo.), 324 (dépo.), 382 (dépo.)\par
\hangindent=1em 17 août : 152 (dépo.), 351 (dépo.)\par
\hangindent=1em 18 août : 157 (indét.), 396 (dépo.)\par
\hangindent=1em 19 août : 137 (dépo.)\par
\hangindent=1em 20 août : 133 (dépo.), 453 (indét.)\par
\hangindent=1em 21 août : 140 (dépo.)\par
\hangindent=1em 23 août : 166 (dépo.), 287 (indét.)\par
\hangindent=1em 24 août : 409 (dépo.)\par
\hangindent=1em 26 août : 431 (dépo.)\par
\hangindent=1em 27 août : 17 (dépo.)\par
\hangindent=1em 31 août : 68 (dépo.), 258 (déc.)\par
\hangindent=1em 4 septembre : 419 (indét.)\par
\hangindent=1em 5 septembre : 115 (dépo.), 139 (dépo.)\par
\hangindent=1em 7 septembre : 47 (dépo.)\par
\hangindent=1em 9 septembre : 451 (indét.)\par
\hangindent=1em 11 septembre : 250 (indét.)\par
\hangindent=1em 12 septembre : 16 (dépo.)\par
\hangindent=1em 13 septembre : 347 (indét.)\par
\hangindent=1em 14 septembre : 206 (dépo.)\par
\hangindent=1em 15 septembre : 15 (déc.), 371 (déc.)\par
\hangindent=1em 16 septembre : 146 (dépo.), 418 (indét.)\par
\hangindent=1em 17 septembre : 292 (dépo.), 349 (indét.)\par
\hangindent=1em 19 septembre : 463 (dépo.)\par
\hangindent=1em 20 septembre : 260 (déc.)\par
\hangindent=1em 21 septembre : 378 (dépo.)\par
\hangindent=1em 23 septembre : 100 (dépo.), 395 (dépo.)\par
\hangindent=1em 25 septembre : 84 (dépo.)\par
\hangindent=1em 26 septembre : 117 (dépo.), 231 (dépo.)\par
\hangindent=1em 27 septembre : 55 (dépo.), 191 (dépo.)\par
\hangindent=1em 28 septembre : 154 (dépo.)\par
\hangindent=1em 29 septembre : 78 (dépo.)\par
\hangindent=1em 30 septembre : 223 (déc.)\par
\hangindent=1em 1 octobre : 323 (dépo.)\par
\hangindent=1em 3 octobre : 132 (dépo.)\par
\hangindent=1em 4 octobre : 125 (dépo.), 134 (indét.), 306 (indét.), 352 (dépo.), 368 (dépo.)\par
\hangindent=1em 8 octobre : 60 (dépo.)\par
\hangindent=1em 9 octobre : 363 (dépo.)\par
\hangindent=1em 10 octobre : 97 (dépo.)\par
\hangindent=1em 11 octobre : 41 (dépo.), 388 (dépo.), 448 (dépo.)\par
\hangindent=1em 13 octobre : 355 (dépo.)\par
\hangindent=1em 15 octobre : 142 (dépo.)\par
\hangindent=1em 18 octobre : 313 (dépo.)\par
\hangindent=1em 23 octobre : 119 (déc.), 408 (indét.)\par
\hangindent=1em 25 octobre : 119 (dépo.)\par
\hangindent=1em 29 octobre : 134 (dépo.)\par
\hangindent=1em 1 novembre : 184 (ach.), 190 (indét.)\par
\hangindent=1em 3 novembre : 203 (dépo.)\par
\hangindent=1em 9 novembre : 265 (indét.)\par
\hangindent=1em 11 novembre : 83 (dépo.)\par
\hangindent=1em 13 novembre : 404 (dépo.)\par
\hangindent=1em 15 novembre : 9 (indét.), 194 (déc.)\par
\hangindent=1em 17 novembre : 128 (dépo.)\par
\hangindent=1em 19 novembre : 23 (dépo.), 321 (indét.)\par
\hangindent=1em 20 novembre : 186 (indét.)\par
\hangindent=1em 22 novembre : 66 (indét.)\par
\hangindent=1em 24 novembre : 210 (dépo.), 269 (déc.)\par
\hangindent=1em 26 novembre : 183 (dépo.)\par
\hangindent=1em 27 novembre : 429 (dépo.), 434 (dépo.)\par
\hangindent=1em 28 novembre : 171 (dépo.)\par
\hangindent=1em 29 novembre : 99 (dépo.)\par
\hangindent=1em 2 décembre : 155 (dépo.)\par
\hangindent=1em 4 décembre : 235 (indét.)\par
\hangindent=1em 7 décembre : 126 (indét.)\par
\hangindent=1em 8 décembre : 81 (dépo.), 359 (indét.)\par
\hangindent=1em 11 décembre : 193 (déc.), 271 (déc.)\par
\hangindent=1em 13 décembre : 317 (dépo.), 367 (dépo.)\par
\hangindent=1em 18 décembre : 279 (dépo.), 428 (dépo.)\par
\hangindent=1em 19 décembre : 56 (déc.)\par
\hangindent=1em 23 décembre : 277 (indét.)\par
\hangindent=1em 25 décembre : 353 (dépo.)\par
\hangindent=1em 30 décembre : 244 (dépo.)\par\end{multicols}


\section{Index des âges}

Âges exprimés en années (table des personnes corrigée, chapitre~\ref{ch:personnes}).

% généré par scripts/04_generated.py
\begin{multicols}{3}\footnotesize\setlength{\parindent}{0pt}\raggedright
\hangindent=1em 1 ans : 211*, 231, 299, 303, 409\par
\hangindent=1em 2 ans : 16, 35, 124, 151, 292, 349, 371, 453*\par
\hangindent=1em 3 ans : 298, 301, 402*\par
\hangindent=1em 4 ans : 156, 197, 235, 431*, 463\par
\hangindent=1em 5 ans : 17, 210, 221*, 255, 330*, 366, 404, 410*\par
\hangindent=1em 6 ans : 128, 143\par
\hangindent=1em 7 ans : 28*, 61, 100, 347, 425*, 467*\par
\hangindent=1em 8 ans : 97, 157, 172\par
\hangindent=1em 9 ans : 94, 145, 318*\par
\hangindent=1em 10 ans : 251*, 278*\par
\hangindent=1em 11 ans : 59, 117, 236*, 246*, 389*\par
\hangindent=1em 12 ans : 68, 236*, 259*\par
\hangindent=1em 13 ans : 317*, 375*, 416*\par
\hangindent=1em 14 ans : 83*, 119, 361*\par
\hangindent=1em 15 ans : 200, 260\par
\hangindent=1em 16 ans : 64*, 92*, 244\par
\hangindent=1em 17 ans : 181*\par
\hangindent=1em 18 ans : 166, 297\par
\hangindent=1em 19 ans : 9, 55*, 102, 126, 142*, 183\par
\hangindent=1em 20 ans : 70, 111, 399, 468*\par
\hangindent=1em 21 ans : 81, 99, 421*\par
\hangindent=1em 22 ans : 42, 186*, 411, 428*\par
\hangindent=1em 23 ans : 85*, 182*, 399\par
\hangindent=1em 24 ans : 118, 269\par
\hangindent=1em 25 ans : 45*, 127*, 167*, 288, 296, 304, 322*, 324, 376, 382*\par
\hangindent=1em 26 ans : 10*, 60, 346\par
\hangindent=1em 27 ans : 115, 116*, 381\par
\hangindent=1em 28 ans : 1*, 72, 140, 276, 354*, 365, 381, 448\par
\hangindent=1em 29 ans : 34\par
\hangindent=1em 30 ans : 7, 15, 132, 136, 193, 368\par
\hangindent=1em 31 ans : 184, 351\par
\hangindent=1em 32 ans : 53*, 133, 272*\par
\hangindent=1em 33 ans : 455*\par
\hangindent=1em 34 ans : 49\par
\hangindent=1em 35 ans : 21*, 54, 90, 317*, 419\par
\hangindent=1em 37 ans : 41, 109\par
\hangindent=1em 38 ans : 191, 198*\par
\hangindent=1em 40 ans : 4, 75*, 223, 225, 443\par
\hangindent=1em 41 ans : 337*\par
\hangindent=1em 42 ans : 310, 391\par
\hangindent=1em 45 ans : 158*, 383, 403\par
\hangindent=1em 46 ans : 202, 343*\par
\hangindent=1em 48 ans : 185\par
\hangindent=1em 50 ans : 91, 306, 372, 400\par
\hangindent=1em 52 ans : 88*\par
\hangindent=1em 53 ans : 38, 390*\par
\hangindent=1em 54 ans : 398\par
\hangindent=1em 55 ans : 169, 286*\par
\hangindent=1em 57 ans : 388\par
\hangindent=1em 60 ans : 147, 264*, 280*, 405, 446, 466\par
\hangindent=1em 62 ans : 144\par
\hangindent=1em 63 ans : 461*\par
\hangindent=1em 65 ans : 33, 268*, 336, 374*, 418\par
\hangindent=1em 66 ans : 237*\par
\hangindent=1em 70 ans : 275, 306, 396*\par
\hangindent=1em 72 ans : 146\par
\hangindent=1em 73 ans : 67*\par
\hangindent=1em 75 ans : 29*, 73*\par
\hangindent=1em 77 ans : 188*\par
\hangindent=1em 80 ans : 271, 353, 427*\par\end{multicols}


\section{Index des graphies}

Lettre déclarée gravée et lettre attendue, selon les corrections signalées par l’EDCS (chapitre~\ref{ch:rep_graphies}).

% généré par scripts/04_generated.py
\begin{multicols}{2}\footnotesize\setlength{\parindent}{0pt}\raggedright
\hangindent=1em A pour e gravé : 205\par
\hangindent=1em A pour h gravé : 215\par
\hangindent=1em A pour i gravé : 470\par
\hangindent=1em A pour o gravé : 205\par
\hangindent=1em ATV pour ibus gravé : 193\par
\hangindent=1em B pour V gravé : 55, 90, 97, 128, 416\par
\hangindent=1em B pour i gravé : 287\par
\hangindent=1em B pour r gravé : 308, 371\par
\hangindent=1em B pour v gravé : 4, 6, 22, 24, 26, 34, 40, 47, 56, 63, 72, 83, 89, 90, 91, 106, 112, 113, 118, 119, 120, 124, 144, 172, 184, 194, 207, 208, 215, 217, 234, 235, 261, 272, 281, 301, 313, 332, 347, 352, 365, 366, 374, 402, 407, 408, 417, 440, 441, 444, 468, 475\par
\hangindent=1em BE pour vi gravé : 66, 458\par
\hangindent=1em C pour G gravé : 211\par
\hangindent=1em C pour K gravé : 9, 128, 279\par
\hangindent=1em C pour e gravé : 319\par
\hangindent=1em C pour g gravé : 19, 35, 99, 201, 221, 231, 304, 306, 352, 353, 354, 359, 366\par
\hangindent=1em C pour i gravé : 340\par
\hangindent=1em C pour q gravé : 184\par
\hangindent=1em CS pour x gravé : 353\par
\hangindent=1em D pour r gravé : 269\par
\hangindent=1em D pour t gravé : 182, 319, 402\par
\hangindent=1em DD pour tt gravé : 72\par
\hangindent=1em E pour F gravé : 219\par
\hangindent=1em E pour S gravé : 441\par
\hangindent=1em E pour a gravé : 28, 214, 279, 466\par
\hangindent=1em E pour c gravé : 319\par
\hangindent=1em E pour i gravé : 27, 40, 50, 54, 70, 76, 100, 111, 125, 306, 353, 364, 368, 388, 391, 443\par
\hangindent=1em EM pour o gravé : 133\par
\hangindent=1em F pour e gravé : 473\par
\hangindent=1em I pour L gravé : 318\par
\hangindent=1em I pour ae gravé : 26\par
\hangindent=1em I pour c gravé : 37\par
\hangindent=1em I pour e gravé : 4, 25, 28, 34, 55, 60, 80, 84, 94, 124, 151, 157, 191, 193, 260, 261, 287, 306, 337, 347, 353, 359, 365, 366, 367, 381, 389, 390, 395, 411, 423, 431, 441, 443, 446\par
\hangindent=1em I pour l gravé : 186, 266, 287, 333, 359\par
\hangindent=1em I pour o gravé : 231, 402\par
\hangindent=1em I pour r gravé : 446, 464\par
\hangindent=1em I pour t gravé : 37, 63, 132, 136, 153, 164, 186, 231, 245, 296, 306, 308, 333, 371, 433\par
\hangindent=1em I pour te gravé : 306\par
\hangindent=1em I pour v gravé : 231\par
\hangindent=1em I pour y gravé : 1, 145, 172, 233, 238, 276, 308, 352, 391, 417, 418, 419, 424\par
\hangindent=1em IS pour ae gravé : 381\par
\hangindent=1em IT pour Fl gravé : 7\par
\hangindent=1em IV pour ly gravé : 371\par
\hangindent=1em K pour C gravé : 62\par
\hangindent=1em K pour ri gravé : 371\par
\hangindent=1em M pour n gravé : 121, 150\par
\hangindent=1em N pour m gravé : 23, 36, 46, 119, 126, 184, 217, 275, 365, 419, 424\par
\hangindent=1em O pour a gravé : 448, 449, 451\par
\hangindent=1em O pour d gravé : 287\par
\hangindent=1em O pour i gravé : 156, 275\par
\hangindent=1em O pour q gravé : 111, 231, 388, 418\par
\hangindent=1em O pour u gravé : 143, 299\par
\hangindent=1em OM pour us gravé : 204\par
\hangindent=1em ON pour um gravé : 320\par
\hangindent=1em P pour b gravé : 299\par
\hangindent=1em P pour d gravé : 135\par
\hangindent=1em PR pour ba gravé : 201\par
\hangindent=1em Q pour c gravé : 1, 184, 287\par
\hangindent=1em QV pour c gravé : 59, 404, 405\par
\hangindent=1em R pour s gravé : 406\par
\hangindent=1em R pour v gravé : 34\par
\hangindent=1em RH pour hr gravé : 306\par
\hangindent=1em RO pour s gravé : 156\par
\hangindent=1em S pour r gravé : 404\par
\hangindent=1em S pour rum gravé : 416\par
\hangindent=1em S pour t gravé : 336\par
\hangindent=1em S pour te gravé : 131\par
\hangindent=1em S pour x gravé : 13, 130, 132, 272, 302, 304\par
\hangindent=1em SI pour is gravé : 309\par
\hangindent=1em T pour ch gravé : 223\par
\hangindent=1em T pour d gravé : 7\par
\hangindent=1em TI pour Fl gravé : 473\par
\hangindent=1em TI pour it gravé : 296\par
\hangindent=1em TS pour st gravé : 60\par
\hangindent=1em U pour i gravé : 115\par
\hangindent=1em U pour o gravé : 260, 282\par
\hangindent=1em U pour y gravé : 232\par
\hangindent=1em V pour b gravé : 85, 123, 125, 276, 340, 357, 380, 382, 470\par
\hangindent=1em V pour i gravé : 44, 157\par
\hangindent=1em V pour n gravé : 85, 309\par
\hangindent=1em V pour o gravé : 32, 54, 55, 69, 94, 104, 130, 161, 191, 193, 209, 225, 275, 281, 288, 289, 299, 320, 322, 323, 327, 336, 346, 347, 353, 359, 365, 374, 375, 376, 379, 400, 419, 421, 427, 431, 443, 466, 470, 473, 475, 478\par
\hangindent=1em V pour y gravé : 223, 239, 338, 339, 340\par
\hangindent=1em VA pour ag gravé : 40\par
\hangindent=1em VM pour o gravé : 291, 402\par
\hangindent=1em VX pour XV gravé : 128, 198\par
\hangindent=1em Y pour V gravé : 142\par
\hangindent=1em Y pour e gravé : 198, 303\par
\hangindent=1em Y pour i gravé : 368\par
\hangindent=1em Y pour u gravé : 349\par\end{multicols}


\section{Index des mots grecs}

% généré par scripts/04_generated.py
\begin{multicols}{3}\footnotesize\setlength{\parindent}{0pt}\raggedright
\hangindent=1em αὐτοῦ 447\par
\hangindent=1em βεναινᾶτος 110\par
\hangindent=1em βιταλία 168\par
\hangindent=1em βίου 110\par
\hangindent=1em εἰρήνῃ 110\par
\hangindent=1em εἴκοσι 168, 348\par
\hangindent=1em εὐγενεστάτη 168\par
\hangindent=1em εὐσεβίου 21\par
\hangindent=1em ζέσας 348\par
\hangindent=1em ζήσας 110, 348\par
\hangindent=1em ζῶντος 168\par
\hangindent=1em θεοδοσίῳ 360\par
\hangindent=1em θεοφίλα 110\par
\hangindent=1em θεὸν 447\par
\hangindent=1em θῖνε 447\par
\hangindent=1em καλανδῶν 8\par
\hangindent=1em κατάθεσις 8\par
\hangindent=1em κατάκιτε 168, 348\par
\hangindent=1em καὶ 21, 110, 168\par
\hangindent=1em κομήτων 110\par
\hangindent=1em κρίσει 447\par
\hangindent=1em κωνσταντινουπολιτίσσα 168\par
\hangindent=1em κώμης 348\par
\hangindent=1em κῖτε 110\par
\hangindent=1em λουππικίνου 110\par
\hangindent=1em λέων 168\par
\hangindent=1em μαρτίων 110\par
\hangindent=1em μαΐου 21\par
\hangindent=1em μετὰ 348\par
\hangindent=1em μηνὶ 21, 348\par
\hangindent=1em μικρᾶς 348\par
\hangindent=1em νοενβρίου 348\par
\hangindent=1em νου 168\par
\hangindent=1em οεʹ 110\par
\hangindent=1em παντοκράτορα 447\par
\hangindent=1em πεδὶν 168\par
\hangindent=1em πρὸ 8, 110\par
\hangindent=1em πέντε 21\par
\hangindent=1em σεπτεμβρίων 8\par
\hangindent=1em σκυλῖσε 447\par
\hangindent=1em συμβίῳ 110\par
\hangindent=1em τελευτᾷ 348\par
\hangindent=1em τινὰ 447\par
\hangindent=1em τις 447\par
\hangindent=1em τολμήσει 447\par
\hangindent=1em τοῦ 110\par
\hangindent=1em τριάκοντα 21\par
\hangindent=1em τὸν 447\par
\hangindent=1em τόπος 168\par
\hangindent=1em τόπῳ 110\par
\hangindent=1em τῆς 447\par
\hangindent=1em τῶν 348\par
\hangindent=1em τῷδε 110\par
\hangindent=1em υἱὸς 348\par
\hangindent=1em χωρίῳ 348\par
\hangindent=1em ἀλέξανδρος 348\par
\hangindent=1em ἀμβρωσίου 348\par
\hangindent=1em ἀναστάσεως 447\par
\hangindent=1em ἀνεπάη 110\par
\hangindent=1em ἀπὸ 110\par
\hangindent=1em ἁγίων 348\par
\hangindent=1em ἄλλον 447\par
\hangindent=1em ἐνθάδε 348\par
\hangindent=1em ἐξελτὼν 110\par
\hangindent=1em ἐποίησα 110\par
\hangindent=1em ἐτῶν 21, 168\par
\hangindent=1em ἑπτά 168\par
\hangindent=1em ἔνοχος 447\par
\hangindent=1em ἔστε 447\par
\hangindent=1em ἔτη 110, 348, 360\par
\hangindent=1em ἡλιοδώρα 21\par
\hangindent=1em ἡμαίραν 447\par
\hangindent=1em ἰδῶν 110\par
\hangindent=1em ἰοβίνου 110\par
\hangindent=1em ὀκταβίλλης 8\par
\hangindent=1em ὀστέα 447\par
\hangindent=1em ὁρκίζομε 447\par
\hangindent=1em ὑμᾶς 447\par
\hangindent=1em ὑπατείᾳ 8, 21, 110\par
\hangindent=1em ῥουφίνου 21\par
\hangindent=1em ῥωμουλι 8\par\end{multicols}


\section{Index des mots latins}

Toutes les formes du texte normalisé (\emph{u/v} et \emph{i/j} unifiés), développements et restitutions compris.

% généré par scripts/04_generated.py
\begin{multicols}{3}\footnotesize\setlength{\parindent}{0pt}\raggedright
\hangindent=1em \lat{aauuggustis} 465\par
\hangindent=1em \lat{ab} 85, 133, 478\par
\hangindent=1em \lat{abiit} 244\par
\hangindent=1em \lat{abita} 330\par
\hangindent=1em \lat{abnepos} 378\par
\hangindent=1em \lat{absolutus} 211\par
\hangindent=1em \lat{abstulit} 306, 470\par
\hangindent=1em \lat{abudantia} 230\par
\hangindent=1em \lat{abundanti} 448\par
\hangindent=1em \lat{ac} 62, 83, 115, 266, 431\par
\hangindent=1em \lat{accepit} 232\par
\hangindent=1em \lat{accepti} 249\par
\hangindent=1em \lat{accill} 339\par
\hangindent=1em \lat{accipe} 244, 440\par
\hangindent=1em \lat{accipiunt} 232\par
\hangindent=1em \lat{acerbo} 470\par
\hangindent=1em \lat{acerbum} 414\par
\hangindent=1em \lat{acerbus} 382\par
\hangindent=1em \lat{acindino} 2, 3\par
\hangindent=1em \lat{acindyno} 1\par
\hangindent=1em \lat{acone} 32\par
\hangindent=1em \lat{actu} 65\par
\hangindent=1em \lat{ad} 73, 241, 275, 387, 433\par
\hangindent=1em \lat{addidit} 433\par
\hangindent=1em \lat{adelfius} 36\par
\hangindent=1em \lat{adeoda} 378\par
\hangindent=1em \lat{adeodatae} 297\par
\hangindent=1em \lat{adiit} 352\par
\hangindent=1em \lat{administrante} 321\par
\hangindent=1em \lat{admirabili} 262\par
\hangindent=1em \lat{adque} 155, 268, 269, 352\par
\hangindent=1em \lat{aduenit} 22\par
\hangindent=1em \lat{ae} 284\par
\hangindent=1em \lat{aeide} 267\par
\hangindent=1em \lat{aequiti} 186\par
\hangindent=1em \lat{aet} 241\par
\hangindent=1em \lat{aeterna} 73, 126, 152, 309\par
\hangindent=1em \lat{aeternae} 306, 429\par
\hangindent=1em \lat{aeternalem} 46\par
\hangindent=1em \lat{aeternam} 306\par
\hangindent=1em \lat{aeternum} 440\par
\hangindent=1em \lat{aetheriam} 244\par
\hangindent=1em \lat{aeuo} 133\par
\hangindent=1em \lat{afrodisi} 39\par
\hangindent=1em \lat{agape} 207\par
\hangindent=1em \lat{agapeni} 99\par
\hangindent=1em \lat{agapenis} 134\par
\hangindent=1em \lat{agrige} 147\par
\hangindent=1em \lat{agrippina} 133\par
\hangindent=1em \lat{aici} 402\par
\hangindent=1em \lat{al} 331\par
\hangindent=1em \lat{alba} 340\par
\hangindent=1em \lat{albanensis} 378\par
\hangindent=1em \lat{albini} 15, 16, 17, 18\par
\hangindent=1em \lat{albino} 11, 12, 13, 14, 19\par
\hangindent=1em \lat{alexander} 115\par
\hangindent=1em \lat{alexandra} 333, 404\par
\hangindent=1em \lat{alexandrae} 367\par
\hangindent=1em \lat{alexandriae} 304\par
\hangindent=1em \lat{alicensi} 164\par
\hangindent=1em \lat{alta} 73\par
\hangindent=1em \lat{am} 189, 218\par
\hangindent=1em \lat{amabilis} 472\par
\hangindent=1em \lat{amachius} 23\par
\hangindent=1em \lat{amans} 474\par
\hangindent=1em \lat{amanti} 15, 16, 17, 18\par
\hangindent=1em \lat{amantiaes} 383\par
\hangindent=1em \lat{amantio} 11, 12, 13, 14, 19, 166\par
\hangindent=1em \lat{amaris} 43\par
\hangindent=1em \lat{amator} 4\par
\hangindent=1em \lat{amatori} 155\par
\hangindent=1em \lat{amatrix} 4, 232\par
\hangindent=1em \lat{amicae} 214\par
\hangindent=1em \lat{amicatus} 275\par
\hangindent=1em \lat{amici} 189\par
\hangindent=1em \lat{amicis} 133\par
\hangindent=1em \lat{amicitiae} 85\par
\hangindent=1em \lat{amicitiam} 306\par
\hangindent=1em \lat{amico} 268\par
\hangindent=1em \lat{amicorum} 268, 391\par
\hangindent=1em \lat{amicus} 22, 202, 391\par
\hangindent=1em \lat{amor} 133, 431, 470\par
\hangindent=1em \lat{amore} 244\par
\hangindent=1em \lat{amori} 43\par
\hangindent=1em \lat{amoris} 218, 268\par
\hangindent=1em \lat{ampliato} 399\par
\hangindent=1em \lat{amre} 86\par
\hangindent=1em \lat{anastasa} 197\par
\hangindent=1em \lat{anastasi} 195\par
\hangindent=1em \lat{anastasio} 43, 321\par
\hangindent=1em \lat{ancilla} 214, 306, 458\par
\hangindent=1em \lat{aniciis} 369, 371, 372, 373, 377, 379, 383\par
\hangindent=1em \lat{anicio} 381\par
\hangindent=1em \lat{anicis} 368\par
\hangindent=1em \lat{anicitionis} 37\par
\hangindent=1em \lat{anima} 139, 298\par
\hangindent=1em \lat{animam} 382\par
\hangindent=1em \lat{animi} 133\par
\hangindent=1em \lat{animo} 306\par
\hangindent=1em \lat{annis} 1, 4, 10, 13, 14, 15, 16, 20, 34, 38, 43, 63, 64, 81, 100, 132, 140, 142, 151, 168, 184, 215, 223, 230, 255, 259, 268, 278, 285, 306, 309, 315, 330, 354, 389, 390, 391, 399, 429, 464\par
\hangindent=1em \lat{anno} 95, 211, 433, 442\par
\hangindent=1em \lat{annorum} 7, 40, 97, 106, 120, 246, 250, 371, 381, 404, 411, 416, 458\par
\hangindent=1em \lat{annos} 5, 6, 7, 9, 17, 19, 21, 25, 26, 28, 29, 32, 33, 35, 37, 40, 41, 42, 45, 47, 49, 51, 53, 54, 55, 57, 59, 60, 61, 67, 68, 69, 70, 72, 73, 75, 76, 77, 79, 83, 84, 85, 88, 90, 91, 92, 94, 99, 102, 104, 109, 111, 113, 114, 115, 116, 117, 118, 119, 122, 124, 125, 126, 127, 128, 130, 133, 135, 136, 139, 143, 144, 145, 146, 147, 149, 156, 157, 158, 161, 162, 166, 167, 169, 170, 172, 181, 182, 183, 185, 186, 188, 189, 190, 191, 193, 196, 197, 198, 200, 201, 202, 203, 207, 210, 218, 221, 222, 223, 225, 226, 228, 235, 236, 237, 238, 244, 251, 252, 254, 257, 260, 261, 264, 265, 267, 269, 271, 272, 275, 276, 280, 282, 283, 286, 288, 290, 291, 292, 294, 296, 297, 298, 301, 304, 310, 312, 317, 318, 322, 324, 332, 333, 334, 335, 336, 337, 338, 339, 343, 346, 347, 349, 350, 351, 352, 353, 355, 359, 361, 363, 364, 365, 366, 367, 368, 370, 372, 373, 374, 375, 376, 378, 379, 382, 383, 385, 388, 396, 398, 399, 400, 402, 403, 405, 410, 412, 418, 419, 420, 421, 425, 427, 428, 431, 440, 443, 446, 448, 450, 453, 455, 461, 463, 466, 467, 468, 470, 472, 473, 474\par
\hangindent=1em \lat{anns} 329\par
\hangindent=1em \lat{annum} 43, 231, 299, 303, 409, 433\par
\hangindent=1em \lat{ante} 19\par
\hangindent=1em \lat{antia} 181\par
\hangindent=1em \lat{antonio} 228, 229, 230, 231, 232, 233, 234, 235, 236, 237, 238, 239\par
\hangindent=1em \lat{antus} 107\par
\hangindent=1em \lat{apolinus} 102\par
\hangindent=1em \lat{apostoli} 321\par
\hangindent=1em \lat{apriles} 4, 26, 85, 136, 169, 197, 310, 318, 386\par
\hangindent=1em \lat{aprili} 478\par
\hangindent=1em \lat{apro} 185\par
\hangindent=1em \lat{apronianae} 279\par
\hangindent=1em \lat{aqundino} 55\par
\hangindent=1em \lat{aradius} 26\par
\hangindent=1em \lat{arantheo} 155\par
\hangindent=1em \lat{arbitio} 43\par
\hangindent=1em \lat{arbitione} 41, 42, 44\par
\hangindent=1em \lat{arbitioni} 45\par
\hangindent=1em \lat{arcadi} 277, 279, 431, 434, 460, 463, 464\par
\hangindent=1em \lat{arcadio} 276, 280, 281, 343, 345, 346, 347, 363, 367, 373, 385, 386, 387, 389, 390, 391, 392, 393, 394, 395, 461, 462, 465, 466, 467, 468\par
\hangindent=1em \lat{arenteo} 160\par
\hangindent=1em \lat{arinteo} 159\par
\hangindent=1em \lat{arinthei} 156, 158\par
\hangindent=1em \lat{arintheo} 157, 164, 165\par
\hangindent=1em \lat{arinthio} 161\par
\hangindent=1em \lat{aristo} 308\par
\hangindent=1em \lat{armararii} 365\par
\hangindent=1em \lat{armeniacum} 275\par
\hangindent=1em \lat{aronteo} 162, 163\par
\hangindent=1em \lat{artemia} 137\par
\hangindent=1em \lat{artorius} 23\par
\hangindent=1em \lat{as} 261, 364\par
\hangindent=1em \lat{asello} 100\par
\hangindent=1em \lat{asterius} 95\par
\hangindent=1em \lat{astra} 433\par
\hangindent=1em \lat{astris} 433\par
\hangindent=1em \lat{astuto} 280\par
\hangindent=1em \lat{at} 73, 382\par
\hangindent=1em \lat{atque} 26\par
\hangindent=1em \lat{atra} 470\par
\hangindent=1em \lat{attico} 396, 397, 398, 399, 400, 401, 402, 403, 404, 405, 406, 407, 408, 409, 410, 412, 413, 414, 415, 416\par
\hangindent=1em \lat{aufidia} 449\par
\hangindent=1em \lat{aufinianus} 391\par
\hangindent=1em \lat{augendo} 146\par
\hangindent=1em \lat{auggustis} 5, 136, 137, 190, 206, 213, 351, 373, 391, 393, 394, 461\par
\hangindent=1em \lat{auggustorum} 127, 463\par
\hangindent=1em \lat{augustas} 1, 11, 12, 22, 37, 49, 59, 63, 92, 126, 151, 211, 213, 215, 273, 278, 300, 319, 320, 366, 369, 390, 406, 443, 459\par
\hangindent=1em \lat{augusti} 73, 88, 173, 179, 183, 186, 187, 188, 321, 325, 355, 425, 428, 460\par
\hangindent=1em \lat{augustis} 97, 126, 142, 192, 207, 208\par
\hangindent=1em \lat{augusto} 6, 46, 56, 57, 58, 74, 75, 76, 77, 78, 80, 83, 89, 91, 96, 138, 145, 146, 148, 149, 151, 152, 153, 167, 168, 174, 176, 177, 182, 189, 196, 198, 201, 203, 204, 205, 218, 293, 294, 295, 296, 297, 298, 300, 302, 314, 319, 323, 324, 327, 346, 347, 352, 354, 361, 367, 387, 392, 420, 422, 433, 467, 468, 469, 470, 471, 472, 473, 474, 476, 477\par
\hangindent=1em \lat{aui} 24\par
\hangindent=1em \lat{aulo} 339\par
\hangindent=1em \lat{aur} 370\par
\hangindent=1em \lat{aurea} 382\par
\hangindent=1em \lat{aureli} 19, 350\par
\hangindent=1em \lat{aureliae} 6, 23\par
\hangindent=1em \lat{aurelio} 111, 330, 337, 340\par
\hangindent=1em \lat{aurelius} 6, 19, 229\par
\hangindent=1em \lat{ausonio} 210, 211\par
\hangindent=1em \lat{ausustinae} 84\par
\hangindent=1em \lat{aut} 244\par
\hangindent=1em \lat{autem} 19, 20, 382, 399\par
\hangindent=1em \lat{autrix} 78\par
\hangindent=1em \lat{balbilla} 72\par
\hangindent=1em \lat{balerius} 128\par
\hangindent=1em \lat{barbarus} 478\par
\hangindent=1em \lat{barroniano} 89\par
\hangindent=1em \lat{basili} 378\par
\hangindent=1em \lat{basilica} 340, 445\par
\hangindent=1em \lat{baudone} 276, 281\par
\hangindent=1em \lat{baudonis} 279\par
\hangindent=1em \lat{bautone} 280\par
\hangindent=1em \lat{beatum} 433\par
\hangindent=1em \lat{beatus} 381\par
\hangindent=1em \lat{benata} 407\par
\hangindent=1em \lat{bendi} 161\par
\hangindent=1em \lat{bene} 3, 4, 10, 13, 16, 17, 25, 28, 29, 33, 35, 38, 40, 46, 47, 51, 62, 70, 75, 76, 77, 99, 100, 104, 109, 115, 117, 123, 124, 127, 128, 131, 132, 139, 140, 143, 144, 145, 152, 155, 162, 166, 172, 178, 190, 198, 210, 223, 226, 237, 260, 261, 271, 276, 280, 284, 291, 295, 300, 304, 317, 319, 320, 324, 329, 332, 336, 343, 344, 351, 353, 361, 365, 367, 376, 389, 396, 399, 402, 403, 405, 409, 410, 416, 428, 431, 444, 453, 471\par
\hangindent=1em \lat{benedictus} 268, 441\par
\hangindent=1em \lat{benenato} 110\par
\hangindent=1em \lat{benerio} 402\par
\hangindent=1em \lat{benerosa} 323\par
\hangindent=1em \lat{benignus} 352\par
\hangindent=1em \lat{bibi} 121\par
\hangindent=1em \lat{bictora} 4\par
\hangindent=1em \lat{bina} 401\par
\hangindent=1em \lat{bincentiae} 369\par
\hangindent=1em \lat{bis} 250, 397\par
\hangindent=1em \lat{biscandens} 355\par
\hangindent=1em \lat{bisomum} 24, 31, 130, 141, 150, 217, 275, 281, 414\par
\hangindent=1em \lat{bisomus} 374, 478\par
\hangindent=1em \lat{bo} 472\par
\hangindent=1em \lat{bonae} 92, 472\par
\hangindent=1em \lat{boniau} 391\par
\hangindent=1em \lat{bonifae} 405\par
\hangindent=1em \lat{bonifatia} 84\par
\hangindent=1em \lat{bonifatiae} 140\par
\hangindent=1em \lat{bonitatis} 26, 70, 196, 221, 269\par
\hangindent=1em \lat{bonorum} 333\par
\hangindent=1em \lat{bonoso} 16\par
\hangindent=1em \lat{bres} 12, 312\par
\hangindent=1em \lat{busto} 470\par
\hangindent=1em \lat{ca} 73\par
\hangindent=1em \lat{caelestinae} 399\par
\hangindent=1em \lat{caelestis} 382\par
\hangindent=1em \lat{caeli} 244\par
\hangindent=1em \lat{caelia} 204\par
\hangindent=1em \lat{caelo} 73, 382\par
\hangindent=1em \lat{caesare} 46, 56, 58\par
\hangindent=1em \lat{caesari} 411\par
\hangindent=1em \lat{caesario} 396, 397, 398, 399, 400, 401, 402, 403, 404, 405, 406, 407, 408, 409, 410, 413, 414, 415, 416\par
\hangindent=1em \lat{caessario} 412\par
\hangindent=1em \lat{caiani} 22\par
\hangindent=1em \lat{calchedonia} 179\par
\hangindent=1em \lat{calleroe} 193\par
\hangindent=1em \lat{calliste} 81, 181\par
\hangindent=1em \lat{callistratus} 181\par
\hangindent=1em \lat{cam} 434\par
\hangindent=1em \lat{capita} 382\par
\hangindent=1em \lat{cappadocem} 275\par
\hangindent=1em \lat{capricornus} 287\par
\hangindent=1em \lat{cara} 355, 470\par
\hangindent=1em \lat{carabasinus} 57\par
\hangindent=1em \lat{carissimae} 46, 83, 265\par
\hangindent=1em \lat{caritas} 133\par
\hangindent=1em \lat{carmina} 43, 470\par
\hangindent=1em \lat{carne} 73\par
\hangindent=1em \lat{casam} 442\par
\hangindent=1em \lat{casarius} 19\par
\hangindent=1em \lat{cassia} 15\par
\hangindent=1em \lat{casta} 258\par
\hangindent=1em \lat{castae} 218, 310\par
\hangindent=1em \lat{castissima} 49\par
\hangindent=1em \lat{castissimae} 12, 262\par
\hangindent=1em \lat{castitatem} 474\par
\hangindent=1em \lat{castitatis} 78, 382\par
\hangindent=1em \lat{cat} 45\par
\hangindent=1em \lat{catafroniaes} 134\par
\hangindent=1em \lat{catechumino} 405\par
\hangindent=1em \lat{catholicae} 18, 67, 428\par
\hangindent=1em \lat{catulini} 31, 34\par
\hangindent=1em \lat{catulino} 30, 32, 33\par
\hangindent=1em \lat{causa} 268\par
\hangindent=1em \lat{cclarissimis} 10, 14, 32, 71, 162, 220, 221, 225, 234, 236, 243, 252, 261, 284, 306, 307, 308, 317, 337, 339, 369, 370, 373, 374, 375, 376, 377, 381, 383, 400, 402, 409, 412, 416, 418, 448, 475\par
\hangindent=1em \lat{cclarissimorum} 305, 335, 350, 411, 450\par
\hangindent=1em \lat{cconssulibus} 40, 60\par
\hangindent=1em \lat{cconssulum} 156\par
\hangindent=1em \lat{cconsulibus} 54, 114, 235\par
\hangindent=1em \lat{cecili} 102\par
\hangindent=1em \lat{celerine} 231\par
\hangindent=1em \lat{celerini} 44\par
\hangindent=1em \lat{celerinus} 347\par
\hangindent=1em \lat{celestia} 382\par
\hangindent=1em \lat{celestinus} 297\par
\hangindent=1em \lat{cereale} 49\par
\hangindent=1em \lat{cereali} 50\par
\hangindent=1em \lat{cerealis} 52, 319\par
\hangindent=1em \lat{ceriale} 53\par
\hangindent=1em \lat{cessauit} 70\par
\hangindent=1em \lat{cheretio} 17\par
\hangindent=1em \lat{cheretius} 17\par
\hangindent=1em \lat{christi} 38, 43, 140, 151, 306, 458\par
\hangindent=1em \lat{christianae} 201\par
\hangindent=1em \lat{christo} 311\par
\hangindent=1em \lat{christum} 352\par
\hangindent=1em \lat{ciie} 277\par
\hangindent=1em \lat{cinnamius} 202\par
\hangindent=1em \lat{cito} 43\par
\hangindent=1em \lat{ciuem} 275\par
\hangindent=1em \lat{ciuesqit} 106\par
\hangindent=1em \lat{ciuis} 164, 296, 306\par
\hangindent=1em \lat{claearco} 260\par
\hangindent=1em \lat{clarissima} 15, 179, 306, 330\par
\hangindent=1em \lat{clarissimae} 262, 306\par
\hangindent=1em \lat{clarissimi} 44, 173, 188, 279, 321, 325, 362, 425, 428, 430\par
\hangindent=1em \lat{clarissimiae} 19\par
\hangindent=1em \lat{clarissimis} 260, 340, 344, 371, 396, 397, 458\par
\hangindent=1em \lat{clarissimo} 203, 281, 283, 285, 286, 287, 288, 298, 319, 321, 322, 323, 324, 343, 364, 420, 429, 433, 435, 436, 438, 440, 441, 443, 444, 446, 453, 454, 470, 471, 472, 473, 474, 477, 478\par
\hangindent=1em \lat{clarissimorum} 431\par
\hangindent=1em \lat{clarissimus} 36, 179, 213, 430, 452\par
\hangindent=1em \lat{claro} 73\par
\hangindent=1em \lat{claudi} 73\par
\hangindent=1em \lat{claudio} 62, 68, 79, 228, 229, 236\par
\hangindent=1em \lat{clauditur} 470\par
\hangindent=1em \lat{clearchi} 275\par
\hangindent=1em \lat{clearcho} 262, 263, 264, 265, 266, 269, 270, 271, 272, 274\par
\hangindent=1em \lat{clearco} 273\par
\hangindent=1em \lat{clodius} 36\par
\hangindent=1em \lat{clusit} 306\par
\hangindent=1em \lat{co} 121, 152, 170, 469\par
\hangindent=1em \lat{cognomine} 470\par
\hangindent=1em \lat{cognosce} 244\par
\hangindent=1em \lat{collaudant} 178\par
\hangindent=1em \lat{columna} 321\par
\hangindent=1em \lat{comparasse} 217\par
\hangindent=1em \lat{comparatus} 424\par
\hangindent=1em \lat{comparauerunt} 137, 346\par
\hangindent=1em \lat{comparauit} 281, 406\par
\hangindent=1em \lat{compare} 19, 130, 184, 365\par
\hangindent=1em \lat{comparem} 275\par
\hangindent=1em \lat{compari} 10, 13, 46, 117, 147, 184\par
\hangindent=1em \lat{con} 18, 73, 454\par
\hangindent=1em \lat{concordiaes} 424\par
\hangindent=1em \lat{concordius} 59, 141\par
\hangindent=1em \lat{condita} 306\par
\hangindent=1em \lat{conditus} 19, 43\par
\hangindent=1em \lat{confirmans} 474\par
\hangindent=1em \lat{coniuge} 40, 217, 320, 364, 420\par
\hangindent=1em \lat{coniugi} 1, 26, 28, 83, 115, 184, 262, 282, 288, 304, 310, 319, 324, 378, 399, 455, 472\par
\hangindent=1em \lat{coniugii} 242\par
\hangindent=1em \lat{coniugio} 200\par
\hangindent=1em \lat{coniunctam} 306\par
\hangindent=1em \lat{coniunx} 43, 125\par
\hangindent=1em \lat{coniux} 33, 47, 49, 333, 382\par
\hangindent=1em \lat{conpleuitque} 433\par
\hangindent=1em \lat{cons} 217\par
\hangindent=1em \lat{conscendere} 244\par
\hangindent=1em \lat{conscribunt} 470\par
\hangindent=1em \lat{consiliator} 333\par
\hangindent=1em \lat{conssulatu} 52, 115, 299, 355, 411, 464\par
\hangindent=1em \lat{conssulatum} 261, 275, 357\par
\hangindent=1em \lat{conssule} 295\par
\hangindent=1em \lat{conssulibus} 1, 2, 3, 4, 6, 10, 12, 13, 14, 20, 23, 26, 27, 28, 30, 32, 33, 37, 38, 39, 44, 49, 53, 56, 57, 58, 63, 64, 67, 68, 71, 72, 76, 77, 79, 81, 89, 94, 96, 97, 100, 101, 104, 107, 122, 124, 125, 126, 129, 130, 132, 133, 134, 137, 139, 142, 148, 149, 152, 153, 157, 164, 167, 168, 171, 177, 192, 197, 198, 202, 204, 205, 206, 207, 210, 211, 213, 217, 225, 226, 228, 234, 238, 239, 241, 243, 245, 246, 250, 252, 254, 256, 257, 258, 259, 262, 263, 264, 270, 274, 276, 283, 286, 289, 290, 291, 306, 307, 308, 312, 313, 314, 315, 316, 317, 324, 328, 333, 338, 339, 341, 342, 343, 347, 351, 367, 373, 376, 377, 379, 381, 383, 384, 387, 388, 391, 393, 394, 400, 403, 404, 405, 406, 412, 413, 414, 415, 416, 419, 422, 423, 438, 448, 456, 458, 461, 465, 469, 470, 474, 475, 477\par
\hangindent=1em \lat{constant} 138\par
\hangindent=1em \lat{constante} 5, 6\par
\hangindent=1em \lat{constantia} 43\par
\hangindent=1em \lat{constantiae} 455\par
\hangindent=1em \lat{constantio} 5, 46, 47, 48, 56, 57, 58\par
\hangindent=1em \lat{constantius} 147\par
\hangindent=1em \lat{constatio} 6\par
\hangindent=1em \lat{constitit} 43\par
\hangindent=1em \lat{consul} 43\par
\hangindent=1em \lat{consulato} 13\par
\hangindent=1em \lat{consulatu} 16, 31, 45, 73, 98, 118, 127, 143, 156, 158, 169, 173, 184, 224, 242, 277, 279, 282, 286, 305, 320, 321, 325, 335, 362, 365, 374, 425, 431, 434, 437, 450, 460, 463\par
\hangindent=1em \lat{consulatum} 15, 17, 18, 34, 102, 109, 179, 182, 183, 185, 186, 187, 188, 288, 301, 350, 428\par
\hangindent=1em \lat{consulatus} 392, 424\par
\hangindent=1em \lat{consule} 141, 176, 249, 297, 300, 302, 311, 364, 366, 419, 421, 426, 429, 432, 433, 435, 436, 440, 441, 443, 444, 446, 453, 454, 478\par
\hangindent=1em \lat{consules} 181\par
\hangindent=1em \lat{consulibus} 5, 7, 11, 12, 15, 19, 22, 24, 25, 29, 42, 46, 47, 55, 59, 62, 65, 66, 69, 75, 78, 80, 82, 86, 87, 93, 95, 99, 103, 105, 106, 108, 111, 112, 113, 116, 117, 120, 121, 135, 136, 138, 140, 144, 145, 146, 147, 151, 155, 159, 160, 161, 162, 165, 170, 172, 178, 181, 189, 190, 191, 193, 195, 196, 199, 200, 201, 203, 208, 212, 215, 218, 220, 221, 223, 229, 230, 232, 236, 244, 247, 248, 251, 255, 260, 265, 266, 269, 271, 272, 273, 278, 281, 284, 285, 287, 288, 293, 298, 303, 304, 309, 310, 322, 323, 329, 330, 331, 332, 334, 336, 337, 344, 345, 346, 352, 358, 359, 361, 363, 368, 372, 375, 380, 382, 385, 389, 390, 396, 397, 398, 399, 401, 402, 407, 408, 409, 410, 417, 418, 420, 424, 449, 451, 452, 455, 459, 462, 468, 471, 472, 473\par
\hangindent=1em \lat{consulis} 430, 439\par
\hangindent=1em \lat{consultor} 22\par
\hangindent=1em \lat{consulum} 45, 460\par
\hangindent=1em \lat{contagia} 433\par
\hangindent=1em \lat{contenta} 244\par
\hangindent=1em \lat{continet} 470\par
\hangindent=1em \lat{contra} 125, 285, 439\par
\hangindent=1em \lat{corde} 470\par
\hangindent=1em \lat{cornificiae} 77\par
\hangindent=1em \lat{corpora} 258\par
\hangindent=1em \lat{corpore} 211\par
\hangindent=1em \lat{corpori} 182\par
\hangindent=1em \lat{corporis} 433\par
\hangindent=1em \lat{corpus} 73, 440\par
\hangindent=1em \lat{credens} 352\par
\hangindent=1em \lat{crescens} 125\par
\hangindent=1em \lat{crescentia} 425\par
\hangindent=1em \lat{crescentiaes} 417\par
\hangindent=1em \lat{crescentii} 378\par
\hangindent=1em \lat{crescentina} 440\par
\hangindent=1em \lat{crescit} 22\par
\hangindent=1em \lat{criscent} 331\par
\hangindent=1em \lat{criscentinae} 132\par
\hangindent=1em \lat{crudele} 352\par
\hangindent=1em \lat{cuarcu} 267\par
\hangindent=1em \lat{cui} 470\par
\hangindent=1em \lat{cuius} 133, 400\par
\hangindent=1em \lat{culpa} 268\par
\hangindent=1em \lat{cultrix} 232\par
\hangindent=1em \lat{cum} 4, 19, 21, 32, 40, 42, 43, 53, 56, 73, 82, 104, 114, 130, 133, 137, 146, 184, 193, 213, 216, 217, 218, 219, 223, 254, 257, 265, 275, 286, 291, 297, 306, 320, 324, 330, 364, 365, 373, 378, 399, 402, 420, 454, 470, 474\par
\hangindent=1em \lat{cunctis} 244\par
\hangindent=1em \lat{cuniculi} 183\par
\hangindent=1em \lat{cupiens} 244\par
\hangindent=1em \lat{cupiunt} 232\par
\hangindent=1em \lat{curant} 470\par
\hangindent=1em \lat{curatore} 321\par
\hangindent=1em \lat{dabit} 306\par
\hangindent=1em \lat{dafini} 337\par
\hangindent=1em \lat{dagalaifi} 109\par
\hangindent=1em \lat{dagalaifo} 99, 100, 101, 103, 104, 105\par
\hangindent=1em \lat{dagalaiphi} 98\par
\hangindent=1em \lat{damasus} 244\par
\hangindent=1em \lat{datam} 306\par
\hangindent=1em \lat{datiani} 52\par
\hangindent=1em \lat{datiano} 49, 50, 53\par
\hangindent=1em \lat{ddominis} 5, 6, 57, 58, 87, 97, 126, 129, 130, 136, 142, 150, 170, 190, 192, 197, 207, 208, 213, 218, 293, 331, 351, 353, 367, 389, 391, 400\par
\hangindent=1em \lat{ddominorum} 127, 460, 464\par
\hangindent=1em \lat{de} 183, 211, 214, 269, 309, 421, 472\par
\hangindent=1em \lat{debebit} 73\par
\hangindent=1em \lat{debere} 126\par
\hangindent=1em \lat{debuit} 126\par
\hangindent=1em \lat{decem} 315\par
\hangindent=1em \lat{decembres} 6, 9, 23, 66, 81, 99, 126, 128, 149, 155, 171, 180, 183, 186, 193, 194, 210, 235, 269, 271, 321, 359, 367, 429, 434\par
\hangindent=1em \lat{decembribus} 317\par
\hangindent=1em \lat{decentio} 12, 37, 38, 39, 40\par
\hangindent=1em \lat{decessit} 4, 7, 19, 28, 31, 54, 58, 92, 119, 193, 194, 197, 260, 268, 449, 453\par
\hangindent=1em \lat{declarant} 382\par
\hangindent=1em \lat{decore} 244\par
\hangindent=1em \lat{decus} 352\par
\hangindent=1em \lat{dedit} 470, 475\par
\hangindent=1em \lat{deflent} 214\par
\hangindent=1em \lat{defuncta} 79, 160, 300, 381, 416\par
\hangindent=1em \lat{defunctus} 20, 33, 136, 139, 269, 276, 381, 430, 478\par
\hangindent=1em \lat{dei} 104, 214, 306, 382, 458, 474\par
\hangindent=1em \lat{dem} 414\par
\hangindent=1em \lat{denis} 43\par
\hangindent=1em \lat{denos} 397\par
\hangindent=1em \lat{denudatae} 442\par
\hangindent=1em \lat{deo} 73\par
\hangindent=1em \lat{depo} 173, 316\par
\hangindent=1em \lat{depos} 345\par
\hangindent=1em \lat{deposit} 213, 422\par
\hangindent=1em \lat{deposita} 1, 10, 12, 13, 21, 23, 26, 34, 35, 38, 46, 49, 50, 53, 54, 66, 73, 78, 79, 83, 84, 85, 91, 97, 99, 109, 115, 122, 124, 125, 127, 132, 137, 138, 140, 143, 151, 157, 167, 172, 175, 179, 181, 199, 200, 204, 207, 208, 211, 213, 221, 225, 226, 230, 233, 241, 243, 244, 246, 255, 257, 258, 272, 279, 282, 288, 290, 291, 292, 297, 298, 301, 303, 304, 310, 311, 315, 318, 323, 324, 330, 337, 346, 351, 353, 355, 357, 367, 371, 373, 376, 378, 382, 386, 387, 390, 399, 404, 409, 411, 413, 416, 425, 444, 449, 451, 459, 460, 461, 470\par
\hangindent=1em \lat{depositae} 431\par
\hangindent=1em \lat{depositi} 48\par
\hangindent=1em \lat{depositio} 18, 44, 89, 114, 134, 154, 195, 237, 240, 317, 319, 383, 448\par
\hangindent=1em \lat{deposito} 396\par
\hangindent=1em \lat{depositus} 16, 17, 20, 22, 25, 27, 29, 37, 40, 41, 47, 52, 55, 57, 59, 60, 63, 67, 68, 69, 75, 76, 77, 81, 90, 93, 98, 100, 102, 105, 111, 116, 117, 119, 126, 128, 130, 131, 133, 139, 142, 146, 149, 152, 155, 156, 158, 166, 169, 171, 181, 185, 187, 188, 191, 196, 202, 203, 205, 206, 210, 231, 234, 235, 242, 251, 261, 264, 265, 273, 278, 280, 283, 286, 295, 302, 311, 313, 317, 320, 325, 328, 329, 336, 343, 347, 350, 352, 362, 363, 364, 365, 366, 368, 370, 372, 375, 378, 385, 388, 391, 392, 395, 402, 405, 408, 410, 412, 418, 420, 421, 426, 428, 429, 433, 434, 436, 441, 443, 446, 463, 466, 467, 469\par
\hangindent=1em \lat{depossio} 4, 169, 183\par
\hangindent=1em \lat{des} 378\par
\hangindent=1em \lat{desolauit} 241\par
\hangindent=1em \lat{deum} 73\par
\hangindent=1em \lat{deuotus} 321\par
\hangindent=1em \lat{deus} 475\par
\hangindent=1em \lat{di} 13, 169\par
\hangindent=1em \lat{diaconi} 134, 306, 378\par
\hangindent=1em \lat{dicasti} 73\par
\hangindent=1em \lat{dicta} 470\par
\hangindent=1em \lat{didyme} 179\par
\hangindent=1em \lat{die} 7, 20, 28, 48, 55, 78, 80, 81, 100, 118, 122, 126, 128, 131, 136, 149, 152, 160, 169, 171, 172, 179, 193, 203, 208, 221, 225, 251, 260, 271, 275, 287, 306, 323, 324, 344, 349, 351, 381, 388, 396, 400, 406, 416, 427, 428, 429, 431, 444, 455, 462, 464\par
\hangindent=1em \lat{diebus} 10, 15, 16, 38, 95, 100, 132, 151, 211, 259, 268, 278, 309, 390, 399\par
\hangindent=1em \lat{diem} 22, 118, 306\par
\hangindent=1em \lat{dierum} 404, 416\par
\hangindent=1em \lat{dies} 9, 17, 19, 21, 26, 30, 32, 33, 34, 35, 40, 42, 55, 59, 61, 68, 69, 72, 73, 74, 83, 84, 85, 88, 91, 92, 94, 102, 103, 119, 124, 125, 126, 127, 128, 142, 143, 145, 147, 156, 157, 164, 166, 169, 181, 184, 185, 186, 188, 189, 193, 196, 197, 202, 224, 235, 244, 257, 260, 261, 276, 282, 287, 288, 290, 291, 292, 298, 299, 301, 305, 310, 316, 322, 327, 330, 335, 337, 344, 347, 373, 397, 399, 409, 420, 423, 433, 453, 461, 463, 470\par
\hangindent=1em \lat{digna} 257, 472\par
\hangindent=1em \lat{dignae} 184\par
\hangindent=1em \lat{dilecta} 244\par
\hangindent=1em \lat{dilectae} 43\par
\hangindent=1em \lat{dilectus} 433\par
\hangindent=1em \lat{diocletiano} 64\par
\hangindent=1em \lat{dionisus} 300\par
\hangindent=1em \lat{dis} 72\par
\hangindent=1em \lat{discessit} 223, 309\par
\hangindent=1em \lat{discordia} 115\par
\hangindent=1em \lat{diuo} 91\par
\hangindent=1em \lat{do} 162, 384\par
\hangindent=1em \lat{dolo} 126\par
\hangindent=1em \lat{dolor} 244, 414\par
\hangindent=1em \lat{dolore} 439\par
\hangindent=1em \lat{dolorem} 440, 470\par
\hangindent=1em \lat{doluit} 231\par
\hangindent=1em \lat{dom} 101\par
\hangindent=1em \lat{domina} 449\par
\hangindent=1em \lat{domini} 73, 101, 179, 279, 320, 321, 400, 425, 431\par
\hangindent=1em \lat{dominis} 56, 80, 94, 139, 193, 260, 327, 358, 417, 419, 467\par
\hangindent=1em \lat{domino} 25, 62, 75, 85, 89, 135, 144, 145, 146, 174, 196, 198, 203, 204, 205, 232, 281, 295, 296, 298, 303, 314, 317, 319, 323, 343, 347, 387, 420, 421, 424, 470, 477\par
\hangindent=1em \lat{dominorum} 186\par
\hangindent=1em \lat{dominus} 102\par
\hangindent=1em \lat{domitio} 157\par
\hangindent=1em \lat{domo} 126\par
\hangindent=1em \lat{domum} 46, 214, 241\par
\hangindent=1em \lat{domus} 73, 155, 309\par
\hangindent=1em \lat{donata} 364, 373\par
\hangindent=1em \lat{donatus} 27\par
\hangindent=1em \lat{donia} 246\par
\hangindent=1em \lat{donum} 189\par
\hangindent=1em \lat{dormientes} 179\par
\hangindent=1em \lat{dormit} 73, 400, 464\par
\hangindent=1em \lat{dormitio} 14, 134\par
\hangindent=1em \lat{duceret} 43\par
\hangindent=1em \lat{dulci} 167\par
\hangindent=1em \lat{dulcis} 139, 152, 268, 298, 355, 378, 400, 470\par
\hangindent=1em \lat{dulcissima} 42, 143, 364\par
\hangindent=1em \lat{dulcissimae} 1, 6, 23, 35, 83, 84, 115, 127, 226, 297, 431, 455, 461\par
\hangindent=1em \lat{dulcissimo} 59, 64, 68, 81, 92, 100, 133, 186, 196, 347, 368\par
\hangindent=1em \lat{dulciti} 134\par
\hangindent=1em \lat{dum} 216, 268\par
\hangindent=1em \lat{duo} 81, 124, 173, 411, 468\par
\hangindent=1em \lat{duobus} 151, 349\par
\hangindent=1em \lat{duodecim} 193\par
\hangindent=1em \lat{duodecimum} 43\par
\hangindent=1em \lat{duos} 19, 35\par
\hangindent=1em \lat{dus} 364\par
\hangindent=1em \lat{ea} 449\par
\hangindent=1em \lat{ebosi} 427\par
\hangindent=1em \lat{ecclesiae} 18, 67, 428\par
\hangindent=1em \lat{ecim} 220\par
\hangindent=1em \lat{ed} 373\par
\hangindent=1em \lat{egnatia} 330\par
\hangindent=1em \lat{ego} 147, 217, 340\par
\hangindent=1em \lat{ei} 345\par
\hangindent=1em \lat{eius} 110, 132, 309, 314, 333, 378, 382, 428\par
\hangindent=1em \lat{eleuteriae} 73\par
\hangindent=1em \lat{elio} 277\par
\hangindent=1em \lat{eloquens} 74\par
\hangindent=1em \lat{elpidius} 206\par
\hangindent=1em \lat{emerita} 339\par
\hangindent=1em \lat{emerunt} 387, 417\par
\hangindent=1em \lat{emit} 141, 209, 340, 400\par
\hangindent=1em \lat{emptum} 411\par
\hangindent=1em \lat{emptus} 478\par
\hangindent=1em \lat{en} 159\par
\hangindent=1em \lat{enim} 152, 268\par
\hangindent=1em \lat{entia} 355\par
\hangindent=1em \lat{eorum} 134, 265\par
\hangindent=1em \lat{epinicenis} 418\par
\hangindent=1em \lat{episcopatu} 387\par
\hangindent=1em \lat{episcopi} 378\par
\hangindent=1em \lat{episcopo} 400\par
\hangindent=1em \lat{episcopus} 321\par
\hangindent=1em \lat{equiti} 173, 179, 183, 184, 187, 188\par
\hangindent=1em \lat{equitii} 185\par
\hangindent=1em \lat{equitio} 172, 174, 177, 178, 181, 182\par
\hangindent=1em \lat{erenes} 78\par
\hangindent=1em \lat{erepta} 244\par
\hangindent=1em \lat{ereptus} 352\par
\hangindent=1em \lat{erit} 306\par
\hangindent=1em \lat{eros} 387\par
\hangindent=1em \lat{es} 414\par
\hangindent=1em \lat{esse} 433\par
\hangindent=1em \lat{esset} 133\par
\hangindent=1em \lat{est} 1, 13, 19, 20, 27, 33, 38, 55, 60, 93, 95, 97, 126, 128, 156, 160, 168, 175, 183, 200, 204, 209, 249, 291, 298, 299, 300, 306, 313, 337, 378, 381, 395, 396, 400, 418, 425, 449, 454, 460, 466, 472\par
\hangindent=1em \lat{esuuiae} 26\par
\hangindent=1em \lat{et} 1, 2, 3, 4, 5, 6, 9, 10, 11, 12, 13, 14, 15, 16, 17, 18, 19, 20, 22, 23, 24, 25, 26, 27, 28, 29, 30, 31, 32, 33, 34, 36, 37, 38, 39, 40, 41, 42, 43, 44, 45, 46, 47, 48, 49, 50, 52, 53, 54, 55, 56, 57, 58, 59, 60, 61, 62, 63, 64, 65, 66, 67, 68, 69, 70, 71, 72, 73, 74, 75, 76, 77, 78, 79, 80, 81, 82, 83, 85, 86, 87, 88, 89, 90, 91, 92, 93, 94, 95, 96, 97, 98, 99, 100, 101, 102, 103, 104, 105, 106, 107, 108, 109, 112, 113, 114, 115, 116, 117, 118, 119, 120, 121, 122, 123, 124, 125, 126, 127, 128, 129, 130, 131, 132, 133, 134, 135, 136, 137, 138, 139, 140, 142, 143, 144, 145, 146, 147, 148, 149, 150, 151, 152, 153, 154, 155, 156, 157, 158, 159, 160, 161, 162, 163, 164, 165, 166, 167, 168, 169, 170, 171, 172, 173, 174, 177, 178, 179, 181, 182, 183, 184, 185, 186, 187, 188, 189, 190, 191, 192, 193, 194, 195, 196, 197, 198, 199, 200, 202, 203, 204, 205, 206, 207, 208, 209, 210, 211, 212, 213, 215, 216, 217, 218, 220, 221, 223, 224, 225, 226, 228, 230, 231, 232, 233, 234, 235, 236, 237, 238, 239, 240, 241, 242, 243, 244, 245, 246, 247, 248, 250, 251, 252, 253, 254, 255, 256, 257, 258, 259, 260, 261, 262, 263, 264, 265, 266, 268, 269, 270, 271, 272, 273, 274, 275, 276, 278, 279, 280, 281, 282, 283, 284, 285, 286, 287, 288, 289, 290, 291, 292, 293, 294, 295, 297, 298, 300, 301, 303, 304, 305, 306, 308, 309, 310, 312, 313, 314, 315, 316, 317, 318, 319, 320, 321, 322, 323, 324, 325, 327, 328, 329, 330, 331, 332, 333, 334, 335, 336, 337, 338, 339, 340, 341, 342, 343, 344, 345, 346, 347, 350, 351, 352, 353, 354, 355, 356, 357, 358, 359, 361, 363, 364, 367, 368, 370, 371, 372, 373, 374, 375, 376, 378, 379, 380, 381, 382, 383, 384, 386, 387, 389, 390, 391, 392, 395, 396, 397, 398, 399, 400, 401, 402, 403, 404, 405, 406, 407, 408, 409, 410, 411, 412, 413, 414, 415, 416, 417, 418, 419, 420, 422, 423, 424, 425, 428, 431, 439, 441, 448, 449, 450, 451, 452, 453, 454, 455, 456, 457, 458, 460, 461, 463, 464, 465, 466, 467, 468, 469, 470, 471, 472, 473, 474, 475, 476, 477\par
\hangindent=1em \lat{etati} 231\par
\hangindent=1em \lat{eu} 302, 342\par
\hangindent=1em \lat{eucaerio} 226\par
\hangindent=1em \lat{euceri} 227\par
\hangindent=1em \lat{eucherii} 224\par
\hangindent=1em \lat{eucherio} 220, 221, 223, 225\par
\hangindent=1em \lat{eudaemon} 62\par
\hangindent=1em \lat{eudoxiae} 265\par
\hangindent=1em \lat{eugeni} 355, 357\par
\hangindent=1em \lat{eugenio} 351, 352, 353, 354, 356, 358, 359, 361\par
\hangindent=1em \lat{euhi} 445\par
\hangindent=1em \lat{euodi} 282\par
\hangindent=1em \lat{euodio} 283, 284, 285, 286, 287, 288, 289, 290\par
\hangindent=1em \lat{euodius} 33\par
\hangindent=1em \lat{eusebi} 309\par
\hangindent=1em \lat{eusebio} 20, 54\par
\hangindent=1em \lat{eusebius} 191\par
\hangindent=1em \lat{eustolius} 213\par
\hangindent=1em \lat{eute} 326\par
\hangindent=1em \lat{euterio} 328\par
\hangindent=1em \lat{euthyciano} 423\par
\hangindent=1em \lat{euticiae} 288\par
\hangindent=1em \lat{eutropio} 291, 292, 293, 294\par
\hangindent=1em \lat{eutyches} 68\par
\hangindent=1em \lat{eutychiani} 425, 428, 431\par
\hangindent=1em \lat{eutychiano} 417, 418, 419, 420, 422, 424\par
\hangindent=1em \lat{eutychius} 352\par
\hangindent=1em \lat{ex} 20, 36, 244, 306\par
\hangindent=1em \lat{exempli} 115, 204\par
\hangindent=1em \lat{eximiae} 269\par
\hangindent=1em \lat{exitus} 22\par
\hangindent=1em \lat{exiuit} 56, 271\par
\hangindent=1em \lat{exspectabatur} 231\par
\hangindent=1em \lat{exsuperantius} 156\par
\hangindent=1em \lat{exsuperia} 49\par
\hangindent=1em \lat{extinctus} 433\par
\hangindent=1em \lat{extricatus} 171\par
\hangindent=1em \lat{exuperantius} 245\par
\hangindent=1em \lat{exutus} 433\par
\hangindent=1em \lat{facere} 126\par
\hangindent=1em \lat{factis} 73\par
\hangindent=1em \lat{falso} 470\par
\hangindent=1em \lat{fama} 22\par
\hangindent=1em \lat{fas} 433, 440\par
\hangindent=1em \lat{fasciola} 421\par
\hangindent=1em \lat{fasciolae} 202\par
\hangindent=1em \lat{fateri} 244\par
\hangindent=1em \lat{fati} 182\par
\hangindent=1em \lat{fatum} 295\par
\hangindent=1em \lat{faustinae} 134\par
\hangindent=1em \lat{faustinus} 205\par
\hangindent=1em \lat{fe} 135\par
\hangindent=1em \lat{februarias} 28, 38, 48, 51, 72, 76, 102, 147, 200, 297, 301, 311, 376, 398, 469\par
\hangindent=1em \lat{februariis} 127, 134, 337\par
\hangindent=1em \lat{fecerunt} 126, 197, 300, 347, 368, 409, 441\par
\hangindent=1em \lat{feci} 46\par
\hangindent=1em \lat{fecisset} 135\par
\hangindent=1em \lat{fecisti} 306\par
\hangindent=1em \lat{fecit} 2, 4, 6, 10, 13, 21, 22, 26, 32, 33, 36, 42, 47, 55, 59, 68, 73, 83, 104, 110, 114, 115, 121, 130, 132, 133, 135, 144, 146, 147, 148, 184, 185, 193, 194, 200, 204, 207, 216, 223, 235, 247, 254, 257, 262, 265, 286, 291, 295, 296, 297, 319, 320, 353, 357, 364, 365, 373, 374, 384, 386, 399, 402, 420, 428, 454\par
\hangindent=1em \lat{feliciae} 351\par
\hangindent=1em \lat{felicior} 22\par
\hangindent=1em \lat{felicis} 102\par
\hangindent=1em \lat{felicissima} 34\par
\hangindent=1em \lat{felicitas} 301, 346, 470\par
\hangindent=1em \lat{felicitatis} 417\par
\hangindent=1em \lat{felix} 22, 81, 133, 135, 229\par
\hangindent=1em \lat{felle} 70\par
\hangindent=1em \lat{femina} 15, 179, 204, 306, 440\par
\hangindent=1em \lat{feminae} 19, 115, 262, 265, 310\par
\hangindent=1em \lat{festo} 117\par
\hangindent=1em \lat{fflauiis} 226\par
\hangindent=1em \lat{ffllauiis} 306\par
\hangindent=1em \lat{ffllauis} 236, 452\par
\hangindent=1em \lat{ffllauuis} 252\par
\hangindent=1em \lat{fide} 70\par
\hangindent=1em \lat{fidei} 133, 306, 391\par
\hangindent=1em \lat{fideli} 1\par
\hangindent=1em \lat{fidelis} 70, 115, 231, 278, 391\par
\hangindent=1em \lat{fidem} 306\par
\hangindent=1em \lat{fides} 382\par
\hangindent=1em \lat{fieri} 268\par
\hangindent=1em \lat{figura} 470\par
\hangindent=1em \lat{fili} 231\par
\hangindent=1em \lat{filia} 42, 84, 213, 306, 378\par
\hangindent=1em \lat{filiae} 6, 13, 35, 99, 127, 143, 221, 297, 309, 378, 399, 431\par
\hangindent=1em \lat{filiis} 259\par
\hangindent=1em \lat{filio} 59, 64, 68, 81, 92, 100, 186, 196, 211, 233, 347, 402, 412, 457\par
\hangindent=1em \lat{filippo} 22, 23, 25, 26, 27, 28, 29, 321\par
\hangindent=1em \lat{filite} 30\par
\hangindent=1em \lat{filius} 268, 353\par
\hangindent=1em \lat{filonetis} 250\par
\hangindent=1em \lat{fine} 470\par
\hangindent=1em \lat{flabiani} 362\par
\hangindent=1em \lat{flaui} 73, 88, 96, 98, 156, 183, 350, 355, 411, 430\par
\hangindent=1em \lat{flauia} 72, 440\par
\hangindent=1em \lat{flauiae} 221\par
\hangindent=1em \lat{flauiani} 365\par
\hangindent=1em \lat{flauiano} 364, 366\par
\hangindent=1em \lat{flauiis} 42, 207, 377, 449\par
\hangindent=1em \lat{flauio} 7, 22, 62, 68, 126, 131, 145, 157, 162, 182, 208, 219, 220, 222, 223, 228, 229, 230, 236, 241, 244, 251, 276, 280, 284, 285, 286, 288, 289, 290, 303, 318, 319, 321, 324, 328, 330, 337, 340, 366, 387, 395, 396, 400, 404, 407, 408, 409, 412, 415, 416, 418, 420, 423, 429, 436, 438, 440, 441, 443, 444, 446, 448, 453, 458, 462, 470, 473\par
\hangindent=1em \lat{flauius} 19, 20, 57, 59, 149, 217\par
\hangindent=1em \lat{fletu} 470\par
\hangindent=1em \lat{fletus} 43, 244\par
\hangindent=1em \lat{florentia} 157, 472\par
\hangindent=1em \lat{florentio} 59, 60, 61\par
\hangindent=1em \lat{flori} 244\par
\hangindent=1em \lat{floridus} 266\par
\hangindent=1em \lat{floruit} 266\par
\hangindent=1em \lat{foedera} 244, 306\par
\hangindent=1em \lat{fortunius} 204\par
\hangindent=1em \lat{fossore} 209, 424, 478\par
\hangindent=1em \lat{fossorem} 217\par
\hangindent=1em \lat{frater} 132, 314\par
\hangindent=1em \lat{fratres} 247, 368, 441\par
\hangindent=1em \lat{fratribus} 457\par
\hangindent=1em \lat{fraudito} 456\par
\hangindent=1em \lat{frauita} 451, 452, 454, 455\par
\hangindent=1em \lat{frauiti} 450\par
\hangindent=1em \lat{frauito} 458\par
\hangindent=1em \lat{frauitta} 448, 449\par
\hangindent=1em \lat{fudit} 231\par
\hangindent=1em \lat{fuerat} 244\par
\hangindent=1em \lat{fuisse} 268\par
\hangindent=1em \lat{fuit} 22, 114, 133, 218, 231, 268, 306, 324\par
\hangindent=1em \lat{funere} 470\par
\hangindent=1em \lat{funus} 352\par
\hangindent=1em \lat{futura} 231\par
\hangindent=1em \lat{gaianus} 20, 94\par
\hangindent=1em \lat{galatea} 151\par
\hangindent=1em \lat{galla} 371\par
\hangindent=1em \lat{gaudentia} 451\par
\hangindent=1em \lat{gaudentius} 34, 310, 372\par
\hangindent=1em \lat{gaudet} 433\par
\hangindent=1em \lat{gelasio} 131\par
\hangindent=1em \lat{geminum} 368\par
\hangindent=1em \lat{generosa} 382\par
\hangindent=1em \lat{genetricis} 244\par
\hangindent=1em \lat{genitor} 470\par
\hangindent=1em \lat{genitoris} 244\par
\hangindent=1em \lat{genus} 378\par
\hangindent=1em \lat{germana} 208\par
\hangindent=1em \lat{germine} 182\par
\hangindent=1em \lat{gerontio} 211\par
\hangindent=1em \lat{gloria} 231\par
\hangindent=1em \lat{grata} 382\par
\hangindent=1em \lat{grataque} 257\par
\hangindent=1em \lat{gratiani} 109, 143, 173, 179, 183, 184, 185, 186, 187, 188\par
\hangindent=1em \lat{gratiano} 99, 100, 101, 103, 104, 105, 144, 145, 146, 147, 148, 149, 150, 151, 152, 153, 154, 172, 174, 175, 176, 177, 178, 181, 182, 195, 196, 197, 198, 199, 200, 201, 202, 203, 204, 205, 213, 215, 216, 217, 218\par
\hangindent=1em \lat{gregori} 134\par
\hangindent=1em \lat{habere} 115\par
\hangindent=1em \lat{habet} 352\par
\hangindent=1em \lat{habetur} 378\par
\hangindent=1em \lat{hac} 126, 309, 310\par
\hangindent=1em \lat{haec} 73, 152, 470\par
\hangindent=1em \lat{haes} 244\par
\hangindent=1em \lat{hanc} 306, 382\par
\hangindent=1em \lat{heraclius} 347\par
\hangindent=1em \lat{herculius} 159\par
\hangindent=1em \lat{heredum} 163\par
\hangindent=1em \lat{hermogenia} 10\par
\hangindent=1em \lat{hermoginiano} 371\par
\hangindent=1em \lat{hesperius} 73\par
\hangindent=1em \lat{heu} 244\par
\hangindent=1em \lat{hic} 10, 22, 60, 78, 95, 151, 204, 209, 213, 214, 244, 255, 257, 269, 298, 299, 303, 318, 333, 352, 374, 388, 390, 396, 400, 421, 428, 429, 433, 450, 460, 466, 467, 470\par
\hangindent=1em \lat{hilara} 116\par
\hangindent=1em \lat{hilarino} 117, 410\par
\hangindent=1em \lat{hilaritati} 324, 399\par
\hangindent=1em \lat{hilaro} 320\par
\hangindent=1em \lat{hispanus} 296\par
\hangindent=1em \lat{hoc} 19, 20, 85, 135, 268, 296, 391\par
\hangindent=1em \lat{homini} 268\par
\hangindent=1em \lat{homo} 391, 474\par
\hangindent=1em \lat{honestae} 388\par
\hangindent=1em \lat{honestissima} 440\par
\hangindent=1em \lat{honnorio} 363\par
\hangindent=1em \lat{honori} 282, 425, 428, 460, 463, 464\par
\hangindent=1em \lat{honorio} 283, 284, 285, 286, 287, 288, 289, 290, 367, 373, 385, 386, 387, 389, 390, 391, 392, 393, 394, 395, 417, 418, 419, 420, 421, 422, 423, 424, 426, 427, 461, 465, 467, 468\par
\hangindent=1em \lat{hora} 287\par
\hangindent=1em \lat{horas} 28, 118, 305\par
\hangindent=1em \lat{hortabatur} 241\par
\hangindent=1em \lat{hospes} 22\par
\hangindent=1em \lat{huic} 306, 391\par
\hangindent=1em \lat{humanis} 268\par
\hangindent=1em \lat{hunc} 126, 148, 214, 433, 454\par
\hangindent=1em \lat{hypatio} 54\par
\hangindent=1em \lat{hyppolytus} 98\par
\hangindent=1em \lat{iacet} 85, 151, 255, 257, 287, 303, 352, 390, 433\par
\hangindent=1em \lat{ianuaras} 58\par
\hangindent=1em \lat{ianuaria} 9, 118\par
\hangindent=1em \lat{ianuariae} 26\par
\hangindent=1em \lat{ianuarias} 24, 33, 56, 187, 188, 244, 277, 279, 291, 298, 350, 353, 425, 428, 446, 466\par
\hangindent=1em \lat{ianuariis} 378\par
\hangindent=1em \lat{ianus} 43\par
\hangindent=1em \lat{iciae} 282\par
\hangindent=1em \lat{id} 113\par
\hangindent=1em \lat{idibus} 70, 211, 317, 347, 378\par
\hangindent=1em \lat{idus} 3, 4, 16, 19, 24, 27, 33, 34, 41, 47, 48, 54, 60, 69, 75, 81, 83, 92, 97, 106, 120, 126, 134, 142, 144, 156, 170, 182, 188, 193, 200, 215, 242, 250, 262, 263, 265, 271, 275, 295, 298, 311, 329, 336, 343, 344, 350, 355, 359, 363, 365, 366, 367, 381, 385, 388, 402, 404, 425, 446, 448, 451, 458, 459, 467, 477\par
\hangindent=1em \lat{ie} 173\par
\hangindent=1em \lat{ies} 317\par
\hangindent=1em \lat{igitur} 391\par
\hangindent=1em \lat{ii} 1, 5, 6, 16, 21, 35, 38, 47, 48, 54, 73, 76, 117, 118, 125, 129, 130, 145, 146, 148, 149, 151, 152, 153, 166, 196, 200, 207, 224, 246, 249, 251, 252, 253, 254, 256, 257, 258, 278, 282, 287, 288, 292, 295, 296, 298, 300, 302, 303, 337, 343, 344, 345, 346, 347, 353, 363, 367, 371, 402, 408, 412, 427, 453, 454\par
\hangindent=1em \lat{iia} 138\par
\hangindent=1em \lat{iidilidce} 153\par
\hangindent=1em \lat{iii} 5, 6, 7, 16, 26, 28, 29, 33, 40, 42, 55, 56, 57, 58, 74, 75, 78, 83, 85, 90, 94, 99, 106, 120, 126, 133, 137, 138, 139, 140, 141, 142, 164, 172, 173, 174, 176, 177, 178, 179, 181, 182, 183, 187, 188, 197, 199, 203, 206, 244, 250, 255, 271, 291, 292, 293, 294, 295, 298, 301, 309, 310, 332, 351, 352, 353, 355, 357, 358, 361, 363, 365, 366, 367, 369, 373, 381, 385, 386, 387, 389, 390, 391, 393, 394, 402, 404, 423, 455, 469\par
\hangindent=1em \lat{iiii} 27, 38, 42, 45, 47, 51, 59, 73, 74, 75, 76, 77, 80, 81, 82, 84, 85, 92, 125, 134, 154, 155, 156, 157, 166, 167, 168, 170, 171, 185, 188, 196, 197, 198, 200, 201, 202, 203, 204, 205, 219, 221, 231, 235, 284, 298, 299, 303, 306, 310, 314, 317, 321, 322, 323, 324, 325, 327, 368, 371, 385, 387, 389, 390, 391, 393, 395, 399, 405, 418, 419, 420, 421, 422, 423, 425, 426, 427, 428, 431, 434, 463, 469\par
\hangindent=1em \lat{iiix} 301\par
\hangindent=1em \lat{iiste} 289\par
\hangindent=1em \lat{iit} 133\par
\hangindent=1em \lat{iix} 66, 180\par
\hangindent=1em \lat{illa} 91\par
\hangindent=1em \lat{imus} 152\par
\hangindent=1em \lat{in} 3, 4, 6, 9, 10, 16, 17, 19, 20, 22, 25, 26, 28, 33, 34, 35, 38, 40, 41, 42, 43, 45, 46, 47, 49, 53, 54, 55, 56, 57, 59, 61, 63, 65, 66, 68, 69, 70, 72, 73, 74, 75, 76, 77, 78, 80, 81, 83, 85, 90, 93, 94, 98, 99, 102, 104, 106, 115, 117, 118, 119, 120, 121, 124, 125, 126, 127, 128, 130, 132, 133, 134, 139, 140, 143, 144, 145, 150, 151, 155, 156, 157, 158, 162, 166, 168, 172, 179, 182, 183, 184, 185, 186, 187, 190, 191, 196, 198, 200, 202, 203, 205, 207, 210, 211, 213, 221, 223, 226, 231, 232, 233, 234, 237, 240, 241, 246, 250, 257, 258, 260, 261, 262, 268, 271, 272, 276, 279, 282, 283, 284, 287, 288, 291, 295, 296, 298, 300, 302, 303, 304, 305, 311, 320, 322, 323, 325, 329, 336, 337, 340, 341, 343, 344, 346, 347, 350, 351, 353, 355, 359, 361, 362, 364, 365, 367, 368, 369, 370, 371, 374, 376, 381, 382, 383, 388, 389, 396, 399, 400, 401, 402, 403, 404, 408, 409, 410, 411, 416, 418, 419, 420, 421, 427, 428, 429, 430, 431, 433, 442, 444, 445, 446, 448, 450, 451, 453, 454, 460, 463, 464, 468, 469, 470, 471\par
\hangindent=1em \lat{ina} 290\par
\hangindent=1em \lat{incomparabili} 36\par
\hangindent=1em \lat{incomparabilis} 125\par
\hangindent=1em \lat{infanti} 302\par
\hangindent=1em \lat{infelices} 93, 126\par
\hangindent=1em \lat{infelix} 470\par
\hangindent=1em \lat{inimitabili} 423\par
\hangindent=1em \lat{inn} 241\par
\hangindent=1em \lat{innocens} 95, 231, 298, 400\par
\hangindent=1em \lat{innocenti} 184, 211, 366, 409\par
\hangindent=1em \lat{innocentia} 85\par
\hangindent=1em \lat{innocentiae} 221, 237, 269, 306, 455\par
\hangindent=1em \lat{innocentissimae} 221\par
\hangindent=1em \lat{innocentissimo} 215, 268\par
\hangindent=1em \lat{innocentium} 231\par
\hangindent=1em \lat{innocentius} 67, 179\par
\hangindent=1em \lat{innocuo} 472\par
\hangindent=1em \lat{innox} 181, 308\par
\hangindent=1em \lat{ino} 445\par
\hangindent=1em \lat{inscripsit} 73\par
\hangindent=1em \lat{integrae} 391\par
\hangindent=1em \lat{integritatis} 391, 455\par
\hangindent=1em \lat{intra} 232\par
\hangindent=1em \lat{inus} 469\par
\hangindent=1em \lat{io} 212\par
\hangindent=1em \lat{ioannti} 261\par
\hangindent=1em \lat{iobiani} 88\par
\hangindent=1em \lat{iosimus} 365\par
\hangindent=1em \lat{iouiano} 83, 85, 86, 87, 89, 90, 91, 247\par
\hangindent=1em \lat{iouina} 138\par
\hangindent=1em \lat{iouini} 115\par
\hangindent=1em \lat{iouino} 106, 107, 108, 111, 112, 113, 114, 116, 117, 118, 119, 120, 121, 122, 124\par
\hangindent=1em \lat{ipsa} 73, 216, 306\par
\hangindent=1em \lat{ipse} 22, 126, 152, 244, 352, 357\par
\hangindent=1em \lat{ipsi} 181\par
\hangindent=1em \lat{irm} 148\par
\hangindent=1em \lat{isiu} 332\par
\hangindent=1em \lat{ista} 214\par
\hangindent=1em \lat{isto} 470\par
\hangindent=1em \lat{it} 241\par
\hangindent=1em \lat{ita} 378\par
\hangindent=1em \lat{iterum} 18, 126, 127, 128, 208, 241, 444\par
\hangindent=1em \lat{itur} 241\par
\hangindent=1em \lat{iu} 203, 225\par
\hangindent=1em \lat{iuit} 144, 433\par
\hangindent=1em \lat{iuli} 442\par
\hangindent=1em \lat{iuliae} 124, 221\par
\hangindent=1em \lat{iuliani} 73\par
\hangindent=1em \lat{iuliano} 46, 47, 48, 56, 57, 58, 74, 75, 76, 77, 78, 79, 80, 81, 82\par
\hangindent=1em \lat{iulias} 7, 15, 19, 27, 34, 42, 55, 90, 105, 130, 170, 198, 262, 276, 282, 288, 331, 343, 399, 405, 416, 449, 467, 476\par
\hangindent=1em \lat{iuliis} 70, 95, 211, 470\par
\hangindent=1em \lat{iulius} 67, 116\par
\hangindent=1em \lat{iuncta} 244\par
\hangindent=1em \lat{iuniani} 4\par
\hangindent=1em \lat{iunias} 54, 67, 73, 94, 124, 308, 309, 330, 357, 381, 399, 458\par
\hangindent=1em \lat{iuniis} 50, 97, 122, 460\par
\hangindent=1em \lat{iuniore} 193\par
\hangindent=1em \lat{iuniori} 191\par
\hangindent=1em \lat{iunus} 86\par
\hangindent=1em \lat{iuris} 22\par
\hangindent=1em \lat{ius} 203\par
\hangindent=1em \lat{iustina} 45\par
\hangindent=1em \lat{iustinus} 346, 350\par
\hangindent=1em \lat{iustus} 395\par
\hangindent=1em \lat{iuuenis} 352, 397\par
\hangindent=1em \lat{iuuenta} 352\par
\hangindent=1em \lat{kalendas} 9, 12, 15, 17, 18, 20, 23, 26, 28, 29, 35, 42, 46, 49, 52, 55, 56, 58, 59, 63, 66, 67, 68, 72, 73, 78, 84, 85, 89, 90, 91, 94, 99, 100, 102, 105, 109, 111, 117, 119, 124, 128, 133, 134, 136, 137, 140, 143, 146, 151, 152, 154, 157, 161, 163, 166, 169, 171, 177, 181, 183, 185, 186, 187, 190, 191, 194, 197, 198, 202, 205, 206, 210, 211, 213, 223, 226, 231, 233, 237, 244, 249, 251, 255, 258, 260, 269, 272, 273, 276, 277, 279, 287, 291, 292, 297, 300, 301, 304, 309, 313, 315, 318, 319, 320, 321, 324, 330, 331, 349, 351, 353, 357, 371, 372, 374, 375, 376, 378, 379, 381, 382, 383, 386, 390, 395, 396, 399, 400, 405, 406, 408, 409, 412, 416, 418, 423, 428, 429, 431, 434, 443, 444, 453, 455, 463, 466, 469, 475, 476, 478\par
\hangindent=1em \lat{kalendis} 95, 97, 127, 134, 184, 323, 337, 460, 470\par
\hangindent=1em \lat{kalndas} 441\par
\hangindent=1em \lat{karisiae} 46\par
\hangindent=1em \lat{karissimo} 247\par
\hangindent=1em \lat{laborum} 78\par
\hangindent=1em \lat{lacrimas} 231, 244\par
\hangindent=1em \lat{lacrimis} 433\par
\hangindent=1em \lat{lacu} 183\par
\hangindent=1em \lat{lacurius} 65\par
\hangindent=1em \lat{laesione} 133\par
\hangindent=1em \lat{lampadius} 247\par
\hangindent=1em \lat{laur} 217\par
\hangindent=1em \lat{laurenti} 89, 237, 284\par
\hangindent=1em \lat{laurentia} 233, 257, 303, 312, 339, 432\par
\hangindent=1em \lat{laurentio} 92\par
\hangindent=1em \lat{laurentius} 428, 463\par
\hangindent=1em \lat{laurus} 73\par
\hangindent=1em \lat{lea} 19, 97, 272, 347\par
\hangindent=1em \lat{lector} 67, 202, 269, 421\par
\hangindent=1em \lat{legenti} 470\par
\hangindent=1em \lat{legitima} 47\par
\hangindent=1em \lat{leo} 7, 60, 194, 295, 333, 387\par
\hangindent=1em \lat{leontia} 125, 460\par
\hangindent=1em \lat{leontio} 10\par
\hangindent=1em \lat{leontius} 276\par
\hangindent=1em \lat{leoparda} 203, 349\par
\hangindent=1em \lat{leopardes} 102\par
\hangindent=1em \lat{leopardo} 135\par
\hangindent=1em \lat{leopardus} 269\par
\hangindent=1em \lat{lepus} 430\par
\hangindent=1em \lat{li} 6, 219, 356\par
\hangindent=1em \lat{libera} 172\par
\hangindent=1em \lat{ligitimus} 100\par
\hangindent=1em \lat{lii} 88\par
\hangindent=1em \lat{liii} 38, 162, 390\par
\hangindent=1em \lat{liiii} 398\par
\hangindent=1em \lat{limeni} 31, 34\par
\hangindent=1em \lat{limenio} 30, 32, 33, 378\par
\hangindent=1em \lat{limenius} 378\par
\hangindent=1em \lat{limfiria} 185\par
\hangindent=1em \lat{limina} 232, 352, 387\par
\hangindent=1em \lat{limine} 433\par
\hangindent=1em \lat{loco} 20\par
\hangindent=1em \lat{locum} 73, 130, 137, 141, 148, 150, 217, 320, 340, 346, 347, 365, 368, 387, 400, 407, 411, 417, 454\par
\hangindent=1em \lat{locus} 37, 48, 168, 209, 281, 379, 395, 396, 418, 424, 430, 464, 478\par
\hangindent=1em \lat{lolliani} 45\par
\hangindent=1em \lat{lolliano} 41, 42, 44\par
\hangindent=1em \lat{loquar} 244\par
\hangindent=1em \lat{loum} 216\par
\hangindent=1em \lat{lu} 162, 169, 286, 329, 334\par
\hangindent=1em \lat{lucceia} 306\par
\hangindent=1em \lat{luce} 309, 400\par
\hangindent=1em \lat{lucem} 43, 244, 470\par
\hangindent=1em \lat{lucida} 318\par
\hangindent=1em \lat{lucifero} 268\par
\hangindent=1em \lat{lucio} 433\par
\hangindent=1em \lat{lucis} 352\par
\hangindent=1em \lat{luii} 388\par
\hangindent=1em \lat{luiiii} 142\par
\hangindent=1em \lat{lumina} 470\par
\hangindent=1em \lat{lumine} 73\par
\hangindent=1em \lat{luna} 287, 349\par
\hangindent=1em \lat{lunae} 169, 275, 381\par
\hangindent=1em \lat{lupicini} 115\par
\hangindent=1em \lat{lupicino} 106, 107, 108, 111, 112, 113, 114, 116, 117, 118, 119, 120, 121, 122, 124\par
\hangindent=1em \lat{lx} 147, 264, 280, 405, 446, 464, 466\par
\hangindent=1em \lat{lxii} 144\par
\hangindent=1em \lat{lxiii} 461\par
\hangindent=1em \lat{lxu} 33, 268, 336, 374, 418\par
\hangindent=1em \lat{lxuii} 203\par
\hangindent=1em \lat{lxx} 275, 306, 396, 420, 429\par
\hangindent=1em \lat{lxxii} 146\par
\hangindent=1em \lat{lxxiii} 67\par
\hangindent=1em \lat{lxxu} 29, 73\par
\hangindent=1em \lat{lxxuii} 188\par
\hangindent=1em \lat{lxxx} 271, 353, 427\par
\hangindent=1em \lat{lybrio} 384\par
\hangindent=1em \lat{ma} 18, 205, 241\par
\hangindent=1em \lat{madescant} 470\par
\hangindent=1em \lat{maerentes} 126\par
\hangindent=1em \lat{mages} 368\par
\hangindent=1em \lat{magnentio} 40\par
\hangindent=1em \lat{magno} 295, 296, 297, 300, 302\par
\hangindent=1em \lat{maias} 10, 29, 57, 58, 75, 106, 111, 120, 129, 134, 143, 156, 172, 179, 182, 249, 255, 259, 261, 295, 322, 329, 336, 365, 372, 381, 402, 412, 441, 461, 478\par
\hangindent=1em \lat{maiis} 415\par
\hangindent=1em \lat{mala} 433\par
\hangindent=1em \lat{male} 470\par
\hangindent=1em \lat{malli} 430, 437\par
\hangindent=1em \lat{mallio} 429, 433, 436, 438\par
\hangindent=1em \lat{mamertino} 62, 63, 64, 65, 66, 67, 68, 69, 70, 71, 72\par
\hangindent=1em \lat{mandasti} 73\par
\hangindent=1em \lat{manes} 470\par
\hangindent=1em \lat{manibus} 72\par
\hangindent=1em \lat{marasius} 41\par
\hangindent=1em \lat{marcelli} 378\par
\hangindent=1em \lat{marcellinae} 194\par
\hangindent=1em \lat{marcellino} 4\par
\hangindent=1em \lat{marcellinus} 55\par
\hangindent=1em \lat{marcena} 152\par
\hangindent=1em \lat{marciana} 454\par
\hangindent=1em \lat{marcus} 46, 197\par
\hangindent=1em \lat{mariniano} 59\par
\hangindent=1em \lat{marito} 21, 29, 32, 42, 53, 77, 82, 114, 133, 185, 193, 219, 223, 244, 254, 257, 291, 297, 324, 373, 399\par
\hangindent=1em \lat{maritum} 454\par
\hangindent=1em \lat{maritus} 26, 125, 144, 184, 386\par
\hangindent=1em \lat{martia} 275\par
\hangindent=1em \lat{martias} 20, 46, 89, 104, 109, 131, 134, 144, 185, 202, 303, 383, 400, 452\par
\hangindent=1em \lat{martielis} 429\par
\hangindent=1em \lat{martiniano} 209\par
\hangindent=1em \lat{martinum} 454\par
\hangindent=1em \lat{martis} 7, 118, 122, 128\par
\hangindent=1em \lat{martyrius} 139\par
\hangindent=1em \lat{mater} 81, 102, 382, 470\par
\hangindent=1em \lat{matre} 306\par
\hangindent=1em \lat{matri} 73, 353, 461\par
\hangindent=1em \lat{matrona} 306, 378, 428\par
\hangindent=1em \lat{maura} 130\par
\hangindent=1em \lat{maxima} 419\par
\hangindent=1em \lat{maximae} 144\par
\hangindent=1em \lat{maximi} 411\par
\hangindent=1em \lat{maximiano} 396\par
\hangindent=1em \lat{maximilla} 306\par
\hangindent=1em \lat{maximimasiae} 320\par
\hangindent=1em \lat{maximo} 72, 295, 296, 297, 300, 302, 311\par
\hangindent=1em \lat{me} 217\par
\hangindent=1em \lat{mea} 217\par
\hangindent=1em \lat{meae} 147\par
\hangindent=1em \lat{mecum} 26, 83, 115, 135, 147, 148, 185\par
\hangindent=1em \lat{medicatur} 470\par
\hangindent=1em \lat{medicus} 296\par
\hangindent=1em \lat{melissus} 26\par
\hangindent=1em \lat{melius} 306\par
\hangindent=1em \lat{membra} 470\par
\hangindent=1em \lat{memor} 306\par
\hangindent=1em \lat{memoriae} 19, 92, 262\par
\hangindent=1em \lat{memoriam} 406\par
\hangindent=1em \lat{menbi} 182\par
\hangindent=1em \lat{mense} 330, 478\par
\hangindent=1em \lat{mensem} 68, 257, 399, 409\par
\hangindent=1em \lat{menses} 5, 9, 17, 21, 26, 28, 29, 32, 33, 34, 35, 38, 40, 42, 45, 53, 54, 55, 57, 59, 61, 72, 73, 76, 77, 81, 83, 84, 85, 88, 90, 91, 92, 94, 102, 113, 115, 117, 124, 125, 126, 128, 133, 143, 147, 152, 156, 157, 166, 181, 182, 184, 185, 188, 189, 196, 198, 202, 208, 221, 230, 231, 235, 236, 238, 244, 251, 256, 260, 261, 276, 282, 288, 292, 294, 297, 298, 299, 301, 303, 308, 310, 317, 324, 332, 335, 337, 347, 349, 352, 353, 359, 365, 366, 375, 402, 410, 423, 425, 431, 433, 437, 455, 466\par
\hangindent=1em \lat{mensibus} 1, 10, 16, 95, 100, 132, 151, 211, 255, 259, 268, 278, 309, 354, 399\par
\hangindent=1em \lat{mensium} 97, 205, 371, 378, 404\par
\hangindent=1em \lat{ment} 382\par
\hangindent=1em \lat{mente} 321\par
\hangindent=1em \lat{mentem} 306\par
\hangindent=1em \lat{mentis} 391\par
\hangindent=1em \lat{mercurius} 422\par
\hangindent=1em \lat{merens} 125\par
\hangindent=1em \lat{merente} 271\par
\hangindent=1em \lat{merenti} 3, 4, 10, 13, 16, 17, 25, 28, 29, 33, 35, 38, 40, 46, 47, 51, 62, 70, 75, 76, 77, 99, 100, 104, 109, 115, 117, 121, 123, 124, 127, 128, 131, 132, 139, 140, 143, 144, 145, 152, 155, 162, 166, 172, 178, 190, 198, 210, 223, 226, 237, 260, 261, 276, 280, 284, 291, 295, 300, 304, 317, 319, 320, 324, 329, 332, 336, 343, 344, 351, 361, 365, 367, 376, 389, 396, 399, 402, 403, 409, 410, 416, 428, 431, 444, 453, 471\par
\hangindent=1em \lat{merit} 457\par
\hangindent=1em \lat{merita} 135, 232\par
\hangindent=1em \lat{merite} 353\par
\hangindent=1em \lat{meritis} 73, 306\par
\hangindent=1em \lat{merobaaude} 201\par
\hangindent=1em \lat{merobaude} 195, 196, 197, 198, 199, 200, 202, 203, 204, 205, 236, 241, 243, 244, 245, 246, 247, 248, 249, 250, 251, 252, 253, 254, 256, 257, 258, 259, 261, 278, 298, 303\par
\hangindent=1em \lat{merobaudis} 242, 299\par
\hangindent=1em \lat{mersit} 470\par
\hangindent=1em \lat{merubaude} 255\par
\hangindent=1em \lat{metuisti} 73\par
\hangindent=1em \lat{meum} 440\par
\hangindent=1em \lat{mf} 74\par
\hangindent=1em \lat{mihi} 357, 378, 382, 440\par
\hangindent=1em \lat{mihique} 49\par
\hangindent=1em \lat{minis} 82\par
\hangindent=1em \lat{minus} 67, 73, 75, 83, 106, 115, 120, 127, 144, 146, 156, 167, 178, 183, 193, 201, 210, 222, 223, 225, 226, 250, 257, 271, 275, 276, 280, 294, 296, 304, 317, 337, 351, 361, 376, 383, 389, 396, 400, 403, 405, 418, 419, 428, 429, 443, 446, 453, 455, 472, 474\par
\hangindent=1em \lat{mirae} 26, 70, 93, 221, 269, 306, 455\par
\hangindent=1em \lat{mire} 146\par
\hangindent=1em \lat{miserae} 244\par
\hangindent=1em \lat{misericors} 133, 231\par
\hangindent=1em \lat{modesti} 156, 158\par
\hangindent=1em \lat{modesto} 155, 159, 160, 161, 162, 163, 164, 165\par
\hangindent=1em \lat{mori} 22, 306\par
\hangindent=1em \lat{moribus} 452\par
\hangindent=1em \lat{moritur} 27\par
\hangindent=1em \lat{mors} 306\par
\hangindent=1em \lat{mortis} 352\par
\hangindent=1em \lat{multa} 73\par
\hangindent=1em \lat{multi} 232\par
\hangindent=1em \lat{multis} 244\par
\hangindent=1em \lat{multum} 433\par
\hangindent=1em \lat{mum} 322\par
\hangindent=1em \lat{mundum} 474\par
\hangindent=1em \lat{munus} 182\par
\hangindent=1em \lat{muscula} 151\par
\hangindent=1em \lat{mutauit} 43\par
\hangindent=1em \lat{na} 232\par
\hangindent=1em \lat{nae} 221\par
\hangindent=1em \lat{nam} 306, 470\par
\hangindent=1em \lat{namque} 433\par
\hangindent=1em \lat{nata} 15, 371\par
\hangindent=1em \lat{natale} 321\par
\hangindent=1em \lat{nati} 470\par
\hangindent=1em \lat{natus} 55, 60, 126, 382\par
\hangindent=1em \lat{ne} 305\par
\hangindent=1em \lat{nec} 433, 470\par
\hangindent=1em \lat{neci} 135\par
\hangindent=1em \lat{neofyta} 172\par
\hangindent=1em \lat{neofyto} 145, 210, 233\par
\hangindent=1em \lat{neofytus} 276, 308, 391\par
\hangindent=1em \lat{neoteri} 320, 321\par
\hangindent=1em \lat{neoterio} 314, 319, 322, 323, 324, 327\par
\hangindent=1em \lat{nepos} 378\par
\hangindent=1em \lat{nescit} 22\par
\hangindent=1em \lat{neuetta} 69\par
\hangindent=1em \lat{neuita} 67\par
\hangindent=1em \lat{neuitta} 62, 63, 64, 65, 66, 68, 70, 71, 72\par
\hangindent=1em \lat{neuterii} 325\par
\hangindent=1em \lat{neuterio} 317\par
\hangindent=1em \lat{nic} 267\par
\hangindent=1em \lat{niceforidi} 167\par
\hangindent=1em \lat{nicetius} 296\par
\hangindent=1em \lat{nicomachi} 365\par
\hangindent=1em \lat{nicomacho} 364\par
\hangindent=1em \lat{nicomaci} 362\par
\hangindent=1em \lat{nicomaco} 366\par
\hangindent=1em \lat{nigriniano} 35, 126\par
\hangindent=1em \lat{ninirit} 149\par
\hangindent=1em \lat{nio} 85\par
\hangindent=1em \lat{nioi} 477\par
\hangindent=1em \lat{nis} 374\par
\hangindent=1em \lat{nisi} 240\par
\hangindent=1em \lat{nisus} 311\par
\hangindent=1em \lat{nnostris} 5, 6, 57, 87, 97, 126, 129, 130, 136, 142, 150, 170, 190, 192, 197, 207, 208, 213, 218, 293, 331, 351, 353, 367, 389, 391, 400\par
\hangindent=1em \lat{nnostrorum} 127, 460, 464\par
\hangindent=1em \lat{nnul} 442\par
\hangindent=1em \lat{nnummero} 7\par
\hangindent=1em \lat{nobili} 88\par
\hangindent=1em \lat{nobilissimi} 282\par
\hangindent=1em \lat{nobilissimo} 132, 133, 134, 283, 284, 285, 286, 287, 288, 290\par
\hangindent=1em \lat{nobis} 46, 472\par
\hangindent=1em \lat{noctis} 287\par
\hangindent=1em \lat{nom} 220\par
\hangindent=1em \lat{nomen} 22, 382\par
\hangindent=1em \lat{nomine} 43, 156, 193, 275, 298, 411, 416, 428, 453\par
\hangindent=1em \lat{non} 231\par
\hangindent=1em \lat{nonas} 1, 5, 7, 10, 22, 38, 51, 57, 76, 80, 125, 131, 132, 134, 147, 155, 172, 179, 203, 208, 235, 259, 282, 288, 303, 306, 308, 310, 311, 352, 368, 369, 399, 419, 427, 433, 436, 449\par
\hangindent=1em \lat{nonio} 396, 400, 404, 407, 409, 412, 415, 416\par
\hangindent=1em \lat{nonis} 10, 50, 115, 122, 139, 430\par
\hangindent=1em \lat{nos} 368\par
\hangindent=1em \lat{nostri} 73, 179, 279, 320, 321, 425, 431\par
\hangindent=1em \lat{nostris} 56, 58, 80, 94, 139, 193, 260, 327, 358, 417, 419, 467\par
\hangindent=1em \lat{nostro} 62, 75, 89, 135, 144, 145, 146, 174, 196, 198, 203, 204, 205, 281, 295, 296, 298, 303, 314, 317, 319, 323, 343, 347, 387, 420, 421, 424, 457, 470, 477\par
\hangindent=1em \lat{nostrorum} 186\par
\hangindent=1em \lat{nostrum} 201\par
\hangindent=1em \lat{nota} 114\par
\hangindent=1em \lat{notus} 133\par
\hangindent=1em \lat{nouem} 34, 433\par
\hangindent=1em \lat{nouembres} 83, 119, 134, 190, 203, 265, 272, 313, 332, 404, 408, 444, 475\par
\hangindent=1em \lat{nouembribus} 184\par
\hangindent=1em \lat{nsacp} 180\par
\hangindent=1em \lat{nta} 79\par
\hangindent=1em \lat{ntia} 148\par
\hangindent=1em \lat{numerarius} 429\par
\hangindent=1em \lat{numero} 20, 29, 75, 97, 132, 306, 329, 336, 389\par
\hangindent=1em \lat{numquam} 85, 115\par
\hangindent=1em \lat{nunc} 73\par
\hangindent=1em \lat{nunita} 306\par
\hangindent=1em \lat{nuper} 382\par
\hangindent=1em \lat{nutrit} 73\par
\hangindent=1em \lat{nutriuit} 232\par
\hangindent=1em \lat{ob} 391\par
\hangindent=1em \lat{obiit} 123\par
\hangindent=1em \lat{obsecutus} 93\par
\hangindent=1em \lat{obsequiis} 241\par
\hangindent=1em \lat{octo} 28, 211, 308\par
\hangindent=1em \lat{octobres} 15, 41, 55, 60, 78, 84, 97, 100, 117, 132, 134, 142, 146, 154, 177, 191, 206, 223, 231, 260, 292, 306, 331, 349, 352, 355, 363, 368, 371, 378, 384, 388, 389, 395, 418, 448, 463\par
\hangindent=1em \lat{octobribus} 323\par
\hangindent=1em \lat{octrobres} 125\par
\hangindent=1em \lat{oculis} 244\par
\hangindent=1em \lat{oei} 107\par
\hangindent=1em \lat{oit} 307\par
\hangindent=1em \lat{olibioni} 260\par
\hangindent=1em \lat{olim} 55\par
\hangindent=1em \lat{olybrio} 210, 211, 212, 368, 370, 371, 372, 373, 374, 375, 376, 378, 379, 380, 381, 382, 383\par
\hangindent=1em \lat{omi} 85\par
\hangindent=1em \lat{omni} 65\par
\hangindent=1em \lat{omnia} 85, 214, 257, 306\par
\hangindent=1em \lat{omnibus} 114, 133, 262, 268, 275, 306\par
\hangindent=1em \lat{omnium} 333\par
\hangindent=1em \lat{on} 435\par
\hangindent=1em \lat{onis} 289\par
\hangindent=1em \lat{opas} 202\par
\hangindent=1em \lat{opera} 73, 265\par
\hangindent=1em \lat{operaria} 4\par
\hangindent=1em \lat{opinionis} 391\par
\hangindent=1em \lat{optimique} 306\par
\hangindent=1em \lat{or} 241\par
\hangindent=1em \lat{ora} 214\par
\hangindent=1em \lat{orationibus} 457\par
\hangindent=1em \lat{ore} 462\par
\hangindent=1em \lat{orio} 378\par
\hangindent=1em \lat{ortum} 43, 470\par
\hangindent=1em \lat{orum} 343\par
\hangindent=1em \lat{os} 397\par
\hangindent=1em \lat{ostend} 241\par
\hangindent=1em \lat{ostendis} 306\par
\hangindent=1em \lat{pacactianus} 73\par
\hangindent=1em \lat{pacatiano} 73\par
\hangindent=1em \lat{pace} 3, 4, 6, 9, 10, 16, 17, 20, 22, 26, 28, 33, 34, 35, 38, 40, 41, 42, 45, 46, 47, 49, 53, 54, 55, 56, 57, 59, 61, 63, 66, 68, 69, 70, 72, 73, 74, 75, 76, 77, 78, 80, 81, 83, 90, 93, 94, 98, 99, 102, 104, 106, 115, 117, 119, 120, 121, 124, 125, 126, 127, 128, 130, 132, 133, 134, 139, 140, 143, 144, 145, 150, 151, 152, 155, 156, 157, 158, 162, 165, 166, 168, 172, 179, 182, 183, 184, 185, 186, 187, 190, 191, 196, 198, 200, 202, 203, 205, 207, 210, 211, 213, 221, 223, 226, 231, 233, 234, 237, 240, 241, 246, 250, 258, 260, 261, 268, 271, 272, 276, 279, 282, 283, 284, 287, 288, 291, 295, 296, 298, 300, 302, 303, 304, 305, 320, 322, 323, 325, 329, 336, 337, 341, 343, 344, 346, 347, 350, 351, 353, 355, 359, 361, 362, 364, 365, 367, 368, 369, 370, 371, 374, 376, 381, 383, 388, 389, 396, 399, 401, 402, 403, 404, 408, 409, 410, 411, 416, 419, 420, 421, 427, 428, 429, 430, 431, 444, 445, 446, 448, 450, 451, 453, 460, 463, 464, 469, 470, 471\par
\hangindent=1em \lat{pacem} 275\par
\hangindent=1em \lat{pacis} 257, 400, 418\par
\hangindent=1em \lat{palladio} 344\par
\hangindent=1em \lat{palma} 73, 382\par
\hangindent=1em \lat{palumbo} 70\par
\hangindent=1em \lat{pannonia} 306\par
\hangindent=1em \lat{panthagate} 382\par
\hangindent=1em \lat{paratissimus} 268\par
\hangindent=1em \lat{parentes} 17, 35, 64, 93, 124, 126, 152, 409\par
\hangindent=1em \lat{parenti} 439\par
\hangindent=1em \lat{parentibus} 74, 152\par
\hangindent=1em \lat{parentis} 189\par
\hangindent=1em \lat{parentum} 244\par
\hangindent=1em \lat{pariter} 306\par
\hangindent=1em \lat{parua} 470\par
\hangindent=1em \lat{pasc} 467\par
\hangindent=1em \lat{pascasio} 400\par
\hangindent=1em \lat{pascasius} 235\par
\hangindent=1em \lat{pascentius} 421\par
\hangindent=1em \lat{pater} 2, 22, 68, 241, 296, 297\par
\hangindent=1em \lat{paternae} 268\par
\hangindent=1em \lat{paterniae} 223\par
\hangindent=1em \lat{patre} 378\par
\hangindent=1em \lat{patri} 203\par
\hangindent=1em \lat{patris} 268\par
\hangindent=1em \lat{paula} 160\par
\hangindent=1em \lat{pauli} 218, 321, 387\par
\hangindent=1em \lat{paulinus} 6\par
\hangindent=1em \lat{paulo} 12, 37, 38, 39\par
\hangindent=1em \lat{paulus} 398\par
\hangindent=1em \lat{pauperum} 4, 202, 232\par
\hangindent=1em \lat{pausat} 382\par
\hangindent=1em \lat{pausauit} 458, 464\par
\hangindent=1em \lat{pax} 330\par
\hangindent=1em \lat{pei} 24\par
\hangindent=1em \lat{pelagia} 246\par
\hangindent=1em \lat{pelecrinus} 56\par
\hangindent=1em \lat{per} 22, 231, 257\par
\hangindent=1em \lat{perenn} 241\par
\hangindent=1em \lat{perennis} 231, 306\par
\hangindent=1em \lat{perfectissimus} 59\par
\hangindent=1em \lat{periit} 3\par
\hangindent=1em \lat{perpetuam} 442\par
\hangindent=1em \lat{perseberatia} 299\par
\hangindent=1em \lat{petiliano} 329, 336\par
\hangindent=1em \lat{petri} 154, 387\par
\hangindent=1em \lat{petrio} 68\par
\hangindent=1em \lat{petronia} 416\par
\hangindent=1em \lat{petroniae} 184\par
\hangindent=1em \lat{petroniaes} 183\par
\hangindent=1em \lat{petronio} 144, 145, 149\par
\hangindent=1em \lat{petronius} 374\par
\hangindent=1em \lat{petrus} 197\par
\hangindent=1em \lat{petulum} 439\par
\hangindent=1em \lat{philippo} 24\par
\hangindent=1em \lat{pia} 472\par
\hangindent=1em \lat{piasque} 231\par
\hangindent=1em \lat{picto} 368\par
\hangindent=1em \lat{pictor} 229\par
\hangindent=1em \lat{pientissimi} 247\par
\hangindent=1em \lat{pignera} 470\par
\hangindent=1em \lat{pignusque} 73\par
\hangindent=1em \lat{pila} 445\par
\hangindent=1em \lat{pisonis} 15\par
\hangindent=1em \lat{pitzinnina} 346\par
\hangindent=1em \lat{pius} 275, 352\par
\hangindent=1em \lat{placido} 7, 9\par
\hangindent=1em \lat{plus} 67, 73, 75, 83, 106, 115, 120, 127, 144, 146, 156, 167, 178, 183, 193, 201, 210, 222, 223, 225, 226, 250, 257, 271, 275, 276, 280, 294, 296, 304, 317, 337, 351, 361, 376, 383, 389, 396, 400, 403, 405, 418, 419, 428, 429, 443, 446, 453, 455, 472, 474\par
\hangindent=1em \lat{po} 7\par
\hangindent=1em \lat{porfurius} 204\par
\hangindent=1em \lat{posita} 168, 298, 299, 454\par
\hangindent=1em \lat{positus} 20, 60, 95, 395, 466\par
\hangindent=1em \lat{possedit} 214\par
\hangindent=1em \lat{post} 15, 16, 17, 18, 34, 102, 109, 179, 180, 182, 183, 184, 185, 186, 187, 188, 244, 261, 275, 301, 350, 357, 397, 416, 428, 431, 434\par
\hangindent=1em \lat{postquam} 433\par
\hangindent=1em \lat{posuerunt} 439\par
\hangindent=1em \lat{posuit} 5, 125, 189\par
\hangindent=1em \lat{potens} 306\par
\hangindent=1em \lat{poterit} 306\par
\hangindent=1em \lat{potuit} 115\par
\hangindent=1em \lat{pq} 152\par
\hangindent=1em \lat{praecessit} 6, 266, 275, 368\par
\hangindent=1em \lat{praefectis} 36\par
\hangindent=1em \lat{praefecto} 72, 306\par
\hangindent=1em \lat{praemia} 352\par
\hangindent=1em \lat{praereptus} 43\par
\hangindent=1em \lat{praesidibus} 20\par
\hangindent=1em \lat{praestant} 470\par
\hangindent=1em \lat{praestantissimo} 155\par
\hangindent=1em \lat{praestat} 244\par
\hangindent=1em \lat{praetextatus} 188\par
\hangindent=1em \lat{praetioso} 458\par
\hangindent=1em \lat{praetoriano} 321\par
\hangindent=1em \lat{praetorio} 306\par
\hangindent=1em \lat{preiecticio} 424\par
\hangindent=1em \lat{preniste} 19\par
\hangindent=1em \lat{presbbyteris} 387\par
\hangindent=1em \lat{presbyter} 105, 310, 428\par
\hangindent=1em \lat{presbyteri} 93, 134, 289, 378\par
\hangindent=1em \lat{presbytero} 378\par
\hangindent=1em \lat{pridie} 3, 4, 16, 46, 57, 68, 80, 131, 142, 143, 208, 223, 235, 242, 258, 295, 308, 311, 343, 385, 392, 419, 433, 446, 449, 458\par
\hangindent=1em \lat{prima} 54, 244, 470\par
\hangindent=1em \lat{primaque} 352\par
\hangindent=1em \lat{primi} 10\par
\hangindent=1em \lat{primia} 254\par
\hangindent=1em \lat{primo} 244, 378, 433\par
\hangindent=1em \lat{primum} 22, 89\par
\hangindent=1em \lat{primus} 10\par
\hangindent=1em \lat{princeps} 22\par
\hangindent=1em \lat{priscus} 132\par
\hangindent=1em \lat{priuata} 38\par
\hangindent=1em \lat{pro} 73, 214, 232, 457\par
\hangindent=1em \lat{probi} 143\par
\hangindent=1em \lat{probinianus} 445\par
\hangindent=1em \lat{probino} 4, 368, 370, 371, 372, 373, 374, 375, 376, 378, 379, 380, 381, 382, 383, 384\par
\hangindent=1em \lat{probo} 144, 145, 146, 147, 148, 149, 150, 151, 152, 153, 154\par
\hangindent=1em \lat{proclino} 387\par
\hangindent=1em \lat{proculo} 1, 2, 3, 55\par
\hangindent=1em \lat{prohibet} 244\par
\hangindent=1em \lat{proiectae} 244\par
\hangindent=1em \lat{promoti} 305\par
\hangindent=1em \lat{promoto} 304, 306, 307, 308, 309, 310, 312, 313, 315, 316, 317, 318\par
\hangindent=1em \lat{pronepos} 378\par
\hangindent=1em \lat{properata} 397\par
\hangindent=1em \lat{propositi} 306\par
\hangindent=1em \lat{propositumque} 306\par
\hangindent=1em \lat{proprio} 150\par
\hangindent=1em \lat{propter} 306\par
\hangindent=1em \lat{protulit} 470\par
\hangindent=1em \lat{prudens} 378\par
\hangindent=1em \lat{pudentiana} 269\par
\hangindent=1em \lat{pudica} 70, 472\par
\hangindent=1em \lat{pudicissimae} 115\par
\hangindent=1em \lat{pudicitia} 114\par
\hangindent=1em \lat{pudicitiae} 382\par
\hangindent=1em \lat{pudore} 244\par
\hangindent=1em \lat{puella} 200, 299, 330, 411, 416, 458\par
\hangindent=1em \lat{puer} 142, 156, 308, 347\par
\hangindent=1em \lat{pueri} 88, 282\par
\hangindent=1em \lat{puero} 111, 132, 133, 134, 283, 284, 285, 286, 287, 288, 290\par
\hangindent=1em \lat{pulchra} 244, 298\par
\hangindent=1em \lat{pupus} 433\par
\hangindent=1em \lat{pura} 70\par
\hangindent=1em \lat{qu} 316, 344\par
\hangindent=1em \lat{qua} 73, 115, 442\par
\hangindent=1em \lat{quabus} 241\par
\hangindent=1em \lat{quadratarii} 189\par
\hangindent=1em \lat{quadrisomum} 340, 365\par
\hangindent=1em \lat{quae} 1, 6, 9, 10, 13, 21, 26, 34, 35, 38, 42, 53, 54, 61, 70, 73, 79, 83, 84, 91, 97, 109, 114, 115, 118, 122, 124, 125, 127, 132, 140, 143, 144, 147, 151, 157, 167, 168, 172, 180, 181, 182, 183, 184, 193, 194, 197, 200, 201, 207, 208, 214, 219, 221, 222, 223, 225, 226, 230, 232, 244, 257, 271, 272, 281, 282, 288, 291, 292, 297, 298, 299, 301, 303, 304, 306, 310, 312, 318, 324, 330, 337, 339, 341, 346, 351, 353, 355, 356, 361, 367, 373, 376, 382, 383, 390, 399, 401, 408, 409, 419, 425, 431, 440, 455, 460, 461, 472\par
\hangindent=1em \lat{quaeque} 73\par
\hangindent=1em \lat{quam} 43, 214\par
\hangindent=1em \lat{quanta} 382\par
\hangindent=1em \lat{quartum} 22, 352, 436\par
\hangindent=1em \lat{quater} 43, 392, 417, 424\par
\hangindent=1em \lat{quattuor} 79, 431, 470\par
\hangindent=1em \lat{quem} 73, 209, 241, 400\par
\hangindent=1em \lat{quescit} 90\par
\hangindent=1em \lat{qui} 2, 4, 16, 17, 20, 22, 25, 28, 29, 33, 37, 40, 41, 43, 47, 55, 56, 57, 59, 60, 63, 64, 68, 74, 75, 81, 88, 90, 92, 93, 94, 95, 100, 102, 111, 116, 117, 119, 126, 128, 130, 131, 133, 139, 142, 145, 146, 149, 152, 156, 162, 163, 166, 169, 177, 181, 185, 186, 189, 190, 191, 196, 197, 202, 203, 210, 211, 215, 231, 234, 235, 237, 245, 249, 251, 260, 261, 263, 268, 269, 276, 278, 280, 283, 286, 296, 302, 306, 308, 317, 322, 328, 329, 333, 334, 336, 338, 343, 347, 350, 352, 354, 357, 359, 363, 364, 365, 366, 368, 372, 374, 375, 388, 391, 395, 396, 398, 400, 402, 403, 405, 410, 412, 418, 420, 421, 427, 428, 429, 433, 443, 446, 448, 450, 463, 466, 467, 469, 470, 472, 474\par
\hangindent=1em \lat{quibus} 74\par
\hangindent=1em \lat{quicescuet} 184\par
\hangindent=1em \lat{quid} 244\par
\hangindent=1em \lat{quiescas} 446\par
\hangindent=1em \lat{quiescenti} 405\par
\hangindent=1em \lat{quiesces} 125\par
\hangindent=1em \lat{quiescet} 42, 59, 102\par
\hangindent=1em \lat{quiescis} 73, 231\par
\hangindent=1em \lat{quiescit} 22, 70, 78, 203, 204, 213, 214, 370, 371, 378, 391, 404, 421, 428, 450\par
\hangindent=1em \lat{quieuit} 406\par
\hangindent=1em \lat{quiminati} 210\par
\hangindent=1em \lat{quindecim} 193\par
\hangindent=1em \lat{quinis} 43\par
\hangindent=1em \lat{quinque} 193, 330, 364\par
\hangindent=1em \lat{quinta} 50\par
\hangindent=1em \lat{quinti} 350\par
\hangindent=1em \lat{quintiano} 47\par
\hangindent=1em \lat{quintiliano} 366\par
\hangindent=1em \lat{quintilianus} 474\par
\hangindent=1em \lat{quinto} 339, 340\par
\hangindent=1em \lat{quintum} 470\par
\hangindent=1em \lat{quiriace} 298\par
\hangindent=1em \lat{quirillus} 275\par
\hangindent=1em \lat{quis} 231\par
\hangindent=1em \lat{quo} 133, 470\par
\hangindent=1em \lat{quoad} 133\par
\hangindent=1em \lat{quod} 126, 204, 232, 470\par
\hangindent=1em \lat{quoduuldeus} 112\par
\hangindent=1em \lat{quoduultdeus} 388\par
\hangindent=1em \lat{quriaceti} 319\par
\hangindent=1em \lat{ranns} 441\par
\hangindent=1em \lat{rapetiga} 296\par
\hangindent=1em \lat{rari} 115, 204, 232\par
\hangindent=1em \lat{ratione} 391\par
\hangindent=1em \lat{rcus} 441\par
\hangindent=1em \lat{re} 256, 413\par
\hangindent=1em \lat{rebus} 268\par
\hangindent=1em \lat{recedens} 73\par
\hangindent=1em \lat{recedit} 439\par
\hangindent=1em \lat{recemede} 271\par
\hangindent=1em \lat{receptus} 73\par
\hangindent=1em \lat{recessit} 10, 118, 120, 231, 291, 308, 350\par
\hangindent=1em \lat{recolat} 306\par
\hangindent=1em \lat{recordationis} 268, 306, 388\par
\hangindent=1em \lat{reddis} 306\par
\hangindent=1em \lat{reddit} 15\par
\hangindent=1em \lat{redeat} 470\par
\hangindent=1em \lat{redemptum} 152\par
\hangindent=1em \lat{redemptus} 152\par
\hangindent=1em \lat{refrigerio} 75\par
\hangindent=1em \lat{regione} 183\par
\hangindent=1em \lat{regionis} 378\par
\hangindent=1em \lat{regna} 382\par
\hangindent=1em \lat{religiosa} 306\par
\hangindent=1em \lat{reliquisti} 214\par
\hangindent=1em \lat{reliquit} 182\par
\hangindent=1em \lat{remigia} 428\par
\hangindent=1em \lat{renouas} 440\par
\hangindent=1em \lat{rentis} 189\par
\hangindent=1em \lat{reor} 433\par
\hangindent=1em \lat{replet} 241\par
\hangindent=1em \lat{repletam} 306\par
\hangindent=1em \lat{requiescet} 333, 474\par
\hangindent=1em \lat{requiescit} 10, 145, 374, 388, 429, 460, 464, 467\par
\hangindent=1em \lat{requiescunt} 258\par
\hangindent=1em \lat{requiesicit} 269\par
\hangindent=1em \lat{requieuit} 72, 80\par
\hangindent=1em \lat{requirunt} 214\par
\hangindent=1em \lat{rerinta} 390\par
\hangindent=1em \lat{res} 465\par
\hangindent=1em \lat{respuens} 474\par
\hangindent=1em \lat{retenetur} 133\par
\hangindent=1em \lat{retentat} 470\par
\hangindent=1em \lat{retinet} 244\par
\hangindent=1em \lat{reueccae} 409\par
\hangindent=1em \lat{reuerendo} 203\par
\hangindent=1em \lat{riae} 358\par
\hangindent=1em \lat{ricomede} 260, 270, 274\par
\hangindent=1em \lat{ricomedi} 272\par
\hangindent=1em \lat{ricomere} 262, 263, 264, 265, 266, 269, 273\par
\hangindent=1em \lat{ricomeres} 275\par
\hangindent=1em \lat{rimorido} 478\par
\hangindent=1em \lat{rini} 93\par
\hangindent=1em \lat{ritus} 152\par
\hangindent=1em \lat{ro} 267, 435\par
\hangindent=1em \lat{romae} 19, 429\par
\hangindent=1em \lat{romanillae} 430\par
\hangindent=1em \lat{romano} 145\par
\hangindent=1em \lat{romanus} 22\par
\hangindent=1em \lat{romi} 397\par
\hangindent=1em \lat{romodoro} 476\par
\hangindent=1em \lat{romuliano} 155\par
\hangindent=1em \lat{romulo} 7, 9\par
\hangindent=1em \lat{roscia} 179\par
\hangindent=1em \lat{ru} 462\par
\hangindent=1em \lat{rufina} 300\par
\hangindent=1em \lat{rufinae} 281\par
\hangindent=1em \lat{rufino} 15, 20, 29, 343, 344, 345, 346, 347\par
\hangindent=1em \lat{rufinus} 142\par
\hangindent=1em \lat{rumordio} 470\par
\hangindent=1em \lat{rumorido} 469, 471, 472, 473, 474, 475, 477\par
\hangindent=1em \lat{rus} 241\par
\hangindent=1em \lat{sabbatius} 22, 386\par
\hangindent=1em \lat{sabina} 200\par
\hangindent=1em \lat{sabiniane} 286\par
\hangindent=1em \lat{sabino} 15\par
\hangindent=1em \lat{sacerdot} 289\par
\hangindent=1em \lat{sacrum} 72\par
\hangindent=1em \lat{saecula} 22\par
\hangindent=1em \lat{saeculari} 442\par
\hangindent=1em \lat{saeculo} 125\par
\hangindent=1em \lat{saeculum} 133\par
\hangindent=1em \lat{saepe} 470\par
\hangindent=1em \lat{salia} 23, 24\par
\hangindent=1em \lat{sallea} 22\par
\hangindent=1em \lat{sallia} 25, 26, 27, 28, 29\par
\hangindent=1em \lat{sallusti} 73\par
\hangindent=1em \lat{sallustio} 74, 75, 76, 77, 78, 79, 80, 81, 82\par
\hangindent=1em \lat{salustio} 10\par
\hangindent=1em \lat{sambatiae} 339\par
\hangindent=1em \lat{sanctae} 67, 133, 399\par
\hangindent=1em \lat{sanctis} 73, 330\par
\hangindent=1em \lat{sanctissimae} 310\par
\hangindent=1em \lat{sanctitatis} 26\par
\hangindent=1em \lat{sancto} 211\par
\hangindent=1em \lat{sanctorum} 232\par
\hangindent=1em \lat{sanctum} 249\par
\hangindent=1em \lat{sanctus} 428\par
\hangindent=1em \lat{santulus} 19\par
\hangindent=1em \lat{sapiens} 298, 352\par
\hangindent=1em \lat{sapientiae} 93, 146\par
\hangindent=1em \lat{sarcofago} 19\par
\hangindent=1em \lat{satis} 244\par
\hangindent=1em \lat{satullo} 111\par
\hangindent=1em \lat{saturni} 381, 428, 462\par
\hangindent=1em \lat{saturnini} 242\par
\hangindent=1em \lat{saturnino} 236, 241, 243, 244, 245, 246, 247, 248, 250, 251, 252, 253, 254, 255, 256, 257, 258, 259, 261, 278, 462\par
\hangindent=1em \lat{saturnis} 126\par
\hangindent=1em \lat{satyro} 403\par
\hangindent=1em \lat{sbi} 357\par
\hangindent=1em \lat{scholastico} 469\par
\hangindent=1em \lat{scholasticus} 469\par
\hangindent=1em \lat{scribit} 43\par
\hangindent=1em \lat{scripsi} 217\par
\hangindent=1em \lat{scripsit} 27, 241\par
\hangindent=1em \lat{scripti} 181\par
\hangindent=1em \lat{scriptis} 10, 221\par
\hangindent=1em \lat{se} 24, 55, 108, 137, 141, 194, 197, 204, 207, 235, 281, 300, 334, 342, 346, 347, 357, 374, 402, 407, 417, 441\par
\hangindent=1em \lat{seberane} 124\par
\hangindent=1em \lat{secessit} 250\par
\hangindent=1em \lat{secunda} 445\par
\hangindent=1em \lat{secundae} 70\par
\hangindent=1em \lat{secundi} 48\par
\hangindent=1em \lat{secundianaes} 66\par
\hangindent=1em \lat{secundinus} 136\par
\hangindent=1em \lat{secundos} 48\par
\hangindent=1em \lat{secura} 73\par
\hangindent=1em \lat{sed} 22, 73, 306, 397\par
\hangindent=1em \lat{sedis} 429\par
\hangindent=1em \lat{sem} 241\par
\hangindent=1em \lat{semis} 118\par
\hangindent=1em \lat{semper} 22, 73, 268\par
\hangindent=1em \lat{sensit} 433\par
\hangindent=1em \lat{sensus} 433\par
\hangindent=1em \lat{septem} 28, 359, 365\par
\hangindent=1em \lat{septembres} 5, 16, 17, 30, 35, 47, 68, 69, 80, 91, 133, 137, 140, 152, 157, 161, 166, 167, 208, 226, 233, 234, 245, 250, 252, 258, 263, 287, 304, 324, 351, 382, 396, 409, 419, 423, 431, 451, 453\par
\hangindent=1em \lat{septembribus} 115, 139, 347\par
\hangindent=1em \lat{septembrium} 160\par
\hangindent=1em \lat{septimum} 287, 431\par
\hangindent=1em \lat{sepulcri} 306\par
\hangindent=1em \lat{sepulcrum} 232, 250\par
\hangindent=1em \lat{sepulturae} 306\par
\hangindent=1em \lat{serbulus} 235\par
\hangindent=1em \lat{sergio} 126\par
\hangindent=1em \lat{seruare} 306\par
\hangindent=1em \lat{setema} 387\par
\hangindent=1em \lat{seuera} 17\par
\hangindent=1em \lat{seuerae} 310\par
\hangindent=1em \lat{seuerianus} 400\par
\hangindent=1em \lat{seuerini} 98\par
\hangindent=1em \lat{sex} 189, 237, 433\par
\hangindent=1em \lat{sexaginta} 237\par
\hangindent=1em \lat{sextilius} 343\par
\hangindent=1em \lat{sextum} 359\par
\hangindent=1em \lat{si} 135, 306, 440\par
\hangindent=1em \lat{sia} 143\par
\hangindent=1em \lat{sibi} 23, 24, 26, 36, 124, 130, 137, 144, 184, 194, 197, 216, 281, 306, 310, 334, 346, 347, 406, 417\par
\hangindent=1em \lat{sic} 470\par
\hangindent=1em \lat{siccato} 470\par
\hangindent=1em \lat{signo} 287\par
\hangindent=1em \lat{sil} 335\par
\hangindent=1em \lat{sileam} 244\par
\hangindent=1em \lat{siluana} 56\par
\hangindent=1em \lat{simplicius} 179\par
\hangindent=1em \lat{simul} 440, 454, 470\par
\hangindent=1em \lat{sine} 70, 133, 259, 268, 470\par
\hangindent=1em \lat{sinu} 382\par
\hangindent=1em \lat{sirica} 459\par
\hangindent=1em \lat{sirici} 387\par
\hangindent=1em \lat{siricius} 321\par
\hangindent=1em \lat{sita} 318\par
\hangindent=1em \lat{situm} 352\par
\hangindent=1em \lat{siui} 457\par
\hangindent=1em \lat{so} 352\par
\hangindent=1em \lat{socia} 78\par
\hangindent=1em \lat{socratiaes} 48\par
\hangindent=1em \lat{socrus} 378\par
\hangindent=1em \lat{solacia} 244\par
\hangindent=1em \lat{solaciumque} 214\par
\hangindent=1em \lat{solacius} 133\par
\hangindent=1em \lat{solidis} 411\par
\hangindent=1em \lat{solis} 287\par
\hangindent=1em \lat{solo} 244\par
\hangindent=1em \lat{solum} 201\par
\hangindent=1em \lat{somno} 241, 257, 400, 418\par
\hangindent=1em \lat{soror} 181\par
\hangindent=1em \lat{sorori} 12\par
\hangindent=1em \lat{sororis} 134, 469\par
\hangindent=1em \lat{sosoro} 119\par
\hangindent=1em \lat{soteres} 247\par
\hangindent=1em \lat{sozomene} 97\par
\hangindent=1em \lat{spac} 391\par
\hangindent=1em \lat{sperat} 306\par
\hangindent=1em \lat{spes} 231\par
\hangindent=1em \lat{spiritum} 249\par
\hangindent=1em \lat{spiritus} 73, 400\par
\hangindent=1em \lat{splendori} 73\par
\hangindent=1em \lat{sponsa} 324\par
\hangindent=1em \lat{sponsata} 324\par
\hangindent=1em \lat{ssima} 357\par
\hangindent=1em \lat{statilia} 42\par
\hangindent=1em \lat{statuis} 333\par
\hangindent=1em \lat{staturina} 386\par
\hangindent=1em \lat{stercoria} 108\par
\hangindent=1em \lat{stilichone} 440, 441, 443, 444, 446\par
\hangindent=1em \lat{stilichonis} 439, 442\par
\hangindent=1em \lat{su} 7\par
\hangindent=1em \lat{sua} 4, 19, 40, 56, 130, 214, 286, 306, 364, 365\par
\hangindent=1em \lat{suae} 10, 73, 265, 310, 364\par
\hangindent=1em \lat{sub} 156, 387, 433, 444, 470\par
\hangindent=1em \lat{subolem} 214\par
\hangindent=1em \lat{subolum} 382\par
\hangindent=1em \lat{successae} 35\par
\hangindent=1em \lat{sui} 268\par
\hangindent=1em \lat{suis} 137, 232, 417\par
\hangindent=1em \lat{sumeret} 43, 470\par
\hangindent=1em \lat{summacho} 341\par
\hangindent=1em \lat{sumptu} 150\par
\hangindent=1em \lat{sunt} 382\par
\hangindent=1em \lat{suo} 29, 133, 135, 184, 218, 219, 223, 244, 265, 320, 373, 378, 402, 420, 439\par
\hangindent=1em \lat{suorum} 352\par
\hangindent=1em \lat{supe} 182\par
\hangindent=1em \lat{super} 306, 433\par
\hangindent=1em \lat{superbia} 85\par
\hangindent=1em \lat{superistitem} 214\par
\hangindent=1em \lat{supra} 10, 181, 221\par
\hangindent=1em \lat{susanna} 330\par
\hangindent=1em \lat{susceptus} 400\par
\hangindent=1em \lat{sussanna} 406\par
\hangindent=1em \lat{suum} 45, 454\par
\hangindent=1em \lat{syacrio} 235\par
\hangindent=1em \lat{syagri} 224\par
\hangindent=1em \lat{syagrio} 219, 221, 222, 223, 225, 226, 228, 229, 230, 231, 232, 233, 234, 236, 237, 238, 239, 240, 328\par
\hangindent=1em \lat{sygario} 220\par
\hangindent=1em \lat{symmachi} 335, 350\par
\hangindent=1em \lat{symmacho} 329, 330, 331, 332, 333, 336, 337, 338, 339, 342\par
\hangindent=1em \lat{symmaco} 334, 340\par
\hangindent=1em \lat{ta} 364\par
\hangindent=1em \lat{taeodore} 6\par
\hangindent=1em \lat{tam} 152\par
\hangindent=1em \lat{tanta} 73, 232\par
\hangindent=1em \lat{tante} 382\par
\hangindent=1em \lat{tantum} 458\par
\hangindent=1em \lat{tantus} 133\par
\hangindent=1em \lat{tas} 73\par
\hangindent=1em \lat{tatiani} 335\par
\hangindent=1em \lat{tatiano} 329, 330, 331, 332, 333, 334, 336, 337, 338, 339, 340, 341, 342, 350\par
\hangindent=1em \lat{tauro} 59, 60, 61\par
\hangindent=1em \lat{te} 73, 152, 231, 382, 472\par
\hangindent=1em \lat{tellure} 152\par
\hangindent=1em \lat{tempore} 306, 433\par
\hangindent=1em \lat{temporis} 268\par
\hangindent=1em \lat{tenebris} 43\par
\hangindent=1em \lat{teneros} 241\par
\hangindent=1em \lat{tenet} 382\par
\hangindent=1em \lat{teni} 12\par
\hangindent=1em \lat{ter} 136, 184, 359, 392\par
\hangindent=1em \lat{terio} 24\par
\hangindent=1em \lat{terrae} 73\par
\hangindent=1em \lat{terreni} 433\par
\hangindent=1em \lat{tertium} 147, 193, 416\par
\hangindent=1em \lat{tertulla} 61\par
\hangindent=1em \lat{terus} 287\par
\hangindent=1em \lat{tes} 182\par
\hangindent=1em \lat{thalami} 244\par
\hangindent=1em \lat{theodora} 80, 213\par
\hangindent=1em \lat{theodori} 430, 437\par
\hangindent=1em \lat{theodoro} 429, 432, 433, 436, 438\par
\hangindent=1em \lat{theodosi} 299, 355, 357\par
\hangindent=1em \lat{theodosio} 213, 215, 216, 217, 218, 298, 303, 351, 352, 353, 354, 358, 359, 361, 466, 469, 470, 471, 472, 473, 474, 475, 476, 477\par
\hangindent=1em \lat{theodosius} 452\par
\hangindent=1em \lat{theodotus} 5\par
\hangindent=1em \lat{theofila} 110\par
\hangindent=1em \lat{ti} 287\par
\hangindent=1em \lat{tiae} 180\par
\hangindent=1em \lat{tibi} 330\par
\hangindent=1em \lat{tici} 239\par
\hangindent=1em \lat{ticira} 255\par
\hangindent=1em \lat{tigris} 42\par
\hangindent=1em \lat{tigriti} 353\par
\hangindent=1em \lat{timasi} 305\par
\hangindent=1em \lat{timasio} 304, 306, 308, 310, 312, 313, 315, 316, 317, 318\par
\hangindent=1em \lat{timatio} 309\par
\hangindent=1em \lat{tina} 452\par
\hangindent=1em \lat{titiana} 207\par
\hangindent=1em \lat{titiani} 418\par
\hangindent=1em \lat{tituli} 202\par
\hangindent=1em \lat{tota} 321\par
\hangindent=1em \lat{totius} 155, 237, 382, 391, 455\par
\hangindent=1em \lat{toto} 306\par
\hangindent=1em \lat{tres} 33, 411\par
\hangindent=1em \lat{tribuno} 321\par
\hangindent=1em \lat{triginta} 193\par
\hangindent=1em \lat{trinitatem} 474\par
\hangindent=1em \lat{trisomus} 395\par
\hangindent=1em \lat{tristis} 43\par
\hangindent=1em \lat{trium} 382\par
\hangindent=1em \lat{tu} 382, 440\par
\hangindent=1em \lat{tuae} 231\par
\hangindent=1em \lat{tui} 306\par
\hangindent=1em \lat{tuis} 457\par
\hangindent=1em \lat{tulit} 475\par
\hangindent=1em \lat{tumulata} 440\par
\hangindent=1em \lat{tumulo} 85, 391\par
\hangindent=1em \lat{tumulum} 126, 268\par
\hangindent=1em \lat{tumulus} 244, 470\par
\hangindent=1em \lat{tunc} 306\par
\hangindent=1em \lat{tus} 249, 278\par
\hangindent=1em \lat{tuus} 73\par
\hangindent=1em \lat{uacat} 374\par
\hangindent=1em \lat{ualente} 92, 93, 94, 95, 96, 97, 117, 125, 126, 129, 130, 131, 136, 138, 139, 140, 142, 167, 168, 170, 171, 189, 191, 192, 193, 194, 206, 207, 208, 209\par
\hangindent=1em \lat{ualenti} 128, 166, 169, 190\par
\hangindent=1em \lat{ualentiniani} 102, 127, 169, 320, 321, 325\par
\hangindent=1em \lat{ualentiniano} 92, 93, 94, 95, 96, 97, 117, 125, 126, 128, 129, 130, 131, 132, 133, 134, 135, 136, 137, 138, 139, 140, 141, 142, 166, 167, 168, 170, 171, 189, 190, 191, 192, 193, 194, 206, 207, 208, 209, 291, 292, 293, 294, 301, 314, 317, 319, 322, 323, 324, 327\par
\hangindent=1em \lat{ualentis} 102, 127, 137\par
\hangindent=1em \lat{ualeriae} 109\par
\hangindent=1em \lat{ualerius} 262\par
\hangindent=1em \lat{ualete} 189\par
\hangindent=1em \lat{uarroniani} 88\par
\hangindent=1em \lat{uarroniano} 83, 85, 86, 87, 90\par
\hangindent=1em \lat{uas} 241\par
\hangindent=1em \lat{ubi} 168, 368, 395, 464\par
\hangindent=1em \lat{ubique} 22\par
\hangindent=1em \lat{uda} 273\par
\hangindent=1em \lat{uegetans} 306\par
\hangindent=1em \lat{uel} 85, 189\par
\hangindent=1em \lat{uenerabili} 115\par
\hangindent=1em \lat{uenerandae} 472\par
\hangindent=1em \lat{ueneris} 55, 131, 349, 416, 431\par
\hangindent=1em \lat{ueniam} 306\par
\hangindent=1em \lat{uenit} 19\par
\hangindent=1em \lat{ueram} 305\par
\hangindent=1em \lat{uero} 382\par
\hangindent=1em \lat{uersuum} 382\par
\hangindent=1em \lat{ui} 17, 19, 48, 55, 61, 72, 81, 83, 92, 97, 102, 117, 124, 128, 132, 136, 143, 181, 183, 185, 207, 208, 209, 231, 251, 259, 260, 268, 275, 278, 311, 315, 375, 399, 425, 433, 463, 466\par
\hangindent=1em \lat{uicariae} 429\par
\hangindent=1em \lat{uicentia} 381\par
\hangindent=1em \lat{uictor} 217, 373, 420, 466\par
\hangindent=1em \lat{uictore} 132, 134, 135\par
\hangindent=1em \lat{uictori} 133, 233\par
\hangindent=1em \lat{uictoria} 258, 378\par
\hangindent=1em \lat{uictoriae} 361\par
\hangindent=1em \lat{uictoriaes} 114\par
\hangindent=1em \lat{uictorinae} 19, 127\par
\hangindent=1em \lat{uictorini} 209\par
\hangindent=1em \lat{uictorino} 357\par
\hangindent=1em \lat{uictorio} 412\par
\hangindent=1em \lat{uidebatur} 268\par
\hangindent=1em \lat{uides} 204\par
\hangindent=1em \lat{uiduo} 468\par
\hangindent=1em \lat{uiginti} 34, 70, 411\par
\hangindent=1em \lat{uii} 34, 42, 47, 61, 76, 84, 100, 102, 125, 126, 132, 144, 156, 165, 200, 202, 230, 292, 324, 335, 344, 347, 363, 366, 370, 402, 406, 425, 467\par
\hangindent=1em \lat{uiii} 9, 20, 38, 40, 46, 53, 60, 88, 91, 95, 97, 115, 119, 143, 147, 157, 172, 205, 210, 235, 261, 269, 273, 298, 324, 353, 372, 399, 404, 459, 473\par
\hangindent=1em \lat{uiiii} 20, 48, 58, 94, 100, 111, 117, 119, 145, 166, 183, 202, 300, 318, 344, 365, 395, 399, 409, 412, 420, 433\par
\hangindent=1em \lat{uiiiix} 400\par
\hangindent=1em \lat{uincenti} 450\par
\hangindent=1em \lat{uincentiae} 271\par
\hangindent=1em \lat{uincentio} 448, 449, 451, 452, 453, 456, 458, 459\par
\hangindent=1em \lat{uincentius} 105, 446\par
\hangindent=1em \lat{uinculis} 433\par
\hangindent=1em \lat{uir} 36, 59, 179, 213, 428, 430, 452\par
\hangindent=1em \lat{uirgini} 99\par
\hangindent=1em \lat{uirginia} 4, 286\par
\hangindent=1em \lat{uirginio} 218, 265\par
\hangindent=1em \lat{uirginium} 45\par
\hangindent=1em \lat{uirgo} 70, 183, 306, 385, 425, 458\par
\hangindent=1em \lat{uiri} 22, 44, 173, 188, 279, 306, 321, 325, 362, 425, 428, 430\par
\hangindent=1em \lat{uiricunda} 400\par
\hangindent=1em \lat{uiris} 260, 340, 344, 371, 396, 397, 458\par
\hangindent=1em \lat{uiro} 203, 281, 283, 285, 286, 287, 288, 298, 319, 321, 322, 323, 324, 343, 364, 378, 388, 420, 429, 433, 435, 436, 438, 440, 441, 443, 444, 446, 453, 454, 470, 471, 472, 473, 474, 477\par
\hangindent=1em \lat{uirorum} 431\par
\hangindent=1em \lat{uit} 239\par
\hangindent=1em \lat{uita} 22, 43, 306\par
\hangindent=1em \lat{uitae} 218, 306, 364, 433\par
\hangindent=1em \lat{uitalinae} 376\par
\hangindent=1em \lat{uitalini} 186\par
\hangindent=1em \lat{uitam} 306\par
\hangindent=1em \lat{uiua} 108, 207, 281, 320, 374\par
\hangindent=1em \lat{uiuae} 442\par
\hangindent=1em \lat{uiuenti} 306\par
\hangindent=1em \lat{uiui} 137, 182, 197, 346, 407, 417, 441\par
\hangindent=1em \lat{uiuis} 347\par
\hangindent=1em \lat{uiuit} 22, 73, 306, 469\par
\hangindent=1em \lat{uiuo} 24, 141, 217, 334, 357\par
\hangindent=1em \lat{uiuos} 300\par
\hangindent=1em \lat{uiuum} 402\par
\hangindent=1em \lat{uiuus} 55, 194, 204, 235\par
\hangindent=1em \lat{uix} 268\par
\hangindent=1em \lat{uixi} 329\par
\hangindent=1em \lat{uixit} 1, 2, 4, 5, 6, 7, 9, 10, 13, 14, 15, 16, 17, 19, 20, 21, 25, 26, 28, 29, 33, 34, 35, 37, 38, 40, 41, 42, 45, 47, 49, 53, 54, 55, 57, 59, 60, 61, 63, 64, 67, 68, 70, 72, 73, 74, 75, 76, 79, 81, 83, 84, 85, 88, 90, 91, 92, 93, 94, 95, 99, 100, 102, 109, 111, 115, 116, 117, 118, 119, 122, 124, 125, 126, 127, 128, 130, 131, 132, 133, 136, 139, 140, 142, 143, 144, 145, 146, 147, 149, 151, 156, 157, 158, 161, 162, 163, 166, 167, 168, 169, 170, 172, 181, 182, 183, 184, 185, 186, 188, 189, 190, 191, 193, 194, 196, 197, 198, 200, 201, 202, 203, 207, 208, 210, 211, 215, 219, 221, 222, 223, 225, 226, 228, 230, 231, 234, 235, 236, 237, 244, 245, 251, 252, 255, 257, 259, 260, 261, 263, 264, 269, 271, 272, 275, 276, 278, 280, 282, 283, 286, 288, 291, 292, 296, 297, 298, 299, 301, 302, 303, 304, 305, 306, 308, 310, 312, 317, 318, 322, 324, 328, 330, 332, 333, 334, 335, 336, 337, 338, 339, 343, 346, 347, 349, 350, 351, 353, 354, 355, 356, 359, 361, 363, 364, 365, 366, 367, 368, 370, 372, 373, 374, 375, 376, 378, 382, 383, 385, 388, 389, 390, 391, 396, 398, 399, 400, 401, 402, 403, 405, 409, 410, 412, 418, 419, 420, 421, 423, 425, 427, 428, 429, 431, 433, 440, 442, 443, 446, 448, 450, 453, 455, 460, 461, 463, 466, 467, 468, 469, 472, 473, 474\par
\hangindent=1em \lat{ull} 472\par
\hangindent=1em \lat{um} 163, 182\par
\hangindent=1em \lat{una} 306, 474\par
\hangindent=1em \lat{unae} 442\par
\hangindent=1em \lat{unam} 306\par
\hangindent=1em \lat{uncta} 201\par
\hangindent=1em \lat{uni} 352\par
\hangindent=1em \lat{unius} 196\par
\hangindent=1em \lat{uno} 211, 330\par
\hangindent=1em \lat{untinque} 189\par
\hangindent=1em \lat{unum} 214, 303, 409\par
\hangindent=1em \lat{uocaris} 73\par
\hangindent=1em \lat{uocitus} 133\par
\hangindent=1em \lat{uoluit} 268\par
\hangindent=1em \lat{uolusiana} 449\par
\hangindent=1em \lat{uota} 306\par
\hangindent=1em \lat{uotum} 125, 285, 439\par
\hangindent=1em \lat{ur} 457\par
\hangindent=1em \lat{urania} 241\par
\hangindent=1em \lat{urbem} 22\par
\hangindent=1em \lat{urbi} 72, 93\par
\hangindent=1em \lat{urbicae} 284\par
\hangindent=1em \lat{urbis} 36, 306, 429\par
\hangindent=1em \lat{urgere} 433\par
\hangindent=1em \lat{ursa} 411\par
\hangindent=1em \lat{ursi} 169\par
\hangindent=1em \lat{ursini} 378\par
\hangindent=1em \lat{urso} 387, 478\par
\hangindent=1em \lat{ursus} 124, 184, 443\par
\hangindent=1em \lat{us} 161, 363\par
\hangindent=1em \lat{ut} 93\par
\hangindent=1em \lat{uuiris} 10, 14, 32, 71, 162, 220, 221, 225, 234, 236, 243, 252, 261, 284, 307, 308, 317, 337, 339, 369, 370, 373, 374, 375, 376, 377, 381, 383, 400, 402, 408, 409, 412, 416, 418, 448, 475\par
\hangindent=1em \lat{uuirorum} 305, 335, 350, 411, 450\par
\hangindent=1em \lat{uulnere} 152\par
\hangindent=1em \lat{uuuiris} 306\par
\hangindent=1em \lat{uxor} 10, 34, 110, 400, 438\par
\hangindent=1em \lat{uxore} 56, 104, 146\par
\hangindent=1em \lat{uxori} 23, 36, 167\par
\hangindent=1em \lat{uxoris} 43\par
\hangindent=1em \lat{xciii} 15\par
\hangindent=1em \lat{xi} 33, 35, 49, 55, 57, 59, 64, 77, 89, 97, 105, 117, 128, 188, 224, 236, 242, 246, 297, 299, 309, 324, 349, 376, 378, 383, 389, 411\par
\hangindent=1em \lat{xii} 68, 72, 140, 157, 169, 186, 189, 196, 236, 259, 260, 268, 287, 318, 319, 345, 381, 437, 458, 476\par
\hangindent=1em \lat{xiii} 23, 52, 59, 67, 114, 133, 134, 291, 301, 317, 321, 332, 375, 416, 453, 463\par
\hangindent=1em \lat{xiiii} 9, 12, 56, 77, 83, 109, 119, 137, 147\par
\hangindent=1em \lat{xiu} 361\par
\hangindent=1em \lat{xl} 4, 63, 75, 147, 223, 225, 443\par
\hangindent=1em \lat{xli} 337\par
\hangindent=1em \lat{xlii} 310, 391\par
\hangindent=1em \lat{xliiii} 355\par
\hangindent=1em \lat{xlu} 104, 158, 383, 403\par
\hangindent=1em \lat{xlui} 202, 343\par
\hangindent=1em \lat{xluii} 267\par
\hangindent=1em \lat{xluiii} 185\par
\hangindent=1em \lat{xlx} 372\par
\hangindent=1em \lat{xu} 4, 17, 63, 68, 74, 124, 128, 135, 145, 151, 157, 169, 198, 200, 237, 238, 260, 276, 279, 292, 313, 320, 322, 327, 349, 396, 428, 456, 466\par
\hangindent=1em \lat{xui} 64, 92, 128, 146, 152, 164, 185, 244, 249, 265, 297, 351, 418, 423\par
\hangindent=1em \lat{xuii} 9, 15, 151, 181, 194, 205, 211, 226, 324, 371, 374, 382, 409\par
\hangindent=1em \lat{xuiii} 5, 26, 35, 45, 61, 166, 181, 206, 297, 304, 443, 478\par
\hangindent=1em \lat{xuiiii} 9, 42, 55, 102, 126, 142, 183, 233\par
\hangindent=1em \lat{xx} 16, 30, 103, 111, 178, 185, 197, 223, 257, 260, 305, 399, 461, 468\par
\hangindent=1em \lat{xxi} 19, 40, 81, 99, 100, 181, 399, 421\par
\hangindent=1em \lat{xxii} 19, 42, 83, 186, 399, 428\par
\hangindent=1em \lat{xxiii} 12, 85, 182, 399\par
\hangindent=1em \lat{xxiiii} 118, 269, 278, 290\par
\hangindent=1em \lat{xxu} 13, 26, 45, 69, 114, 127, 167, 168, 201, 244, 282, 288, 296, 304, 322, 324, 376, 382\par
\hangindent=1em \lat{xxui} 10, 37, 60, 94, 126, 127, 309, 346\par
\hangindent=1em \lat{xxuii} 115, 116, 169, 186, 282, 288, 347, 381\par
\hangindent=1em \lat{xxuiii} 1, 72, 125, 140, 250, 276, 354, 365, 381, 399, 448\par
\hangindent=1em \lat{xxuiiii} 352\par
\hangindent=1em \lat{xxx} 7, 15, 132, 136, 146, 199, 259, 287, 368\par
\hangindent=1em \lat{xxxi} 184, 351\par
\hangindent=1em \lat{xxxii} 53, 133, 272\par
\hangindent=1em \lat{xxxiii} 455\par
\hangindent=1em \lat{xxxiiii} 49\par
\hangindent=1em \lat{xxxu} 21, 54, 90, 317, 419, 463\par
\hangindent=1em \lat{xxxuii} 41, 109, 135\par
\hangindent=1em \lat{xxxuiii} 191, 198\par
\hangindent=1em \lat{yperechiam} 275\par
\hangindent=1em \lat{zita} 340\par
\hangindent=1em \lat{zosimi} 424\par
\hangindent=1em \lat{zosimus} 106, 120\par\end{multicols}

